% The next line makes it so it compiles the main file if I hit the compile button
% when editing this file in TeXworks.
% !TEX root = ../AIDMA.tex
%%%%%%%%%%%%%%%%%%%%%%%%%%%%%%%%%%%%
\chapter{Algorithm Analysis}\label{ch:AlgAn}
%%%%%%%%%%%%%%%%%%%%%%%%%%%%%%%%%%%%%%%%%%%%%%%%%%%%%%%%
%This chapter is all my stuff.  CAC, 8/6/13
%Some converted from lecture slides, other stuff from handouts,
% and some stuff that is new.
%%%%%%%%%%%%%%%%%%%%%%%%%%%%%%%%%%%%%%%%%%%%%%%%%%%%%%%%


%%%%%%%%%%%%%%%%%%%%%%%%%%%%%%%%%%%%%%%%%%%%%%%%%%%%%%%%%%%%%%%%%%%%%
In this chapter we take a look at the analysis of algorithms.
The analysis of algorithms is a very important topic in computer science.  
It allows us to determine and express how efficient an algorithm is, 
and it is one of the tools that allows us to compare multiple algorithms
that solve the same problem.

Before we dive into that topic, we first discuss 
one of the most important tools used in algorithm analysis---{\em asymptotic notation}. 
We will define several important notations,
discuss some of the useful properties of the notations, and
provide many examples of two common ways of proving things related to the notations.  
We will then discuss the relative growth rates of several common functions,
focusing on those that are relevant to the topic of algorithm analysis.
We then move on to the most important topic of the chapter in which we 
apply all of this material to the analysis of algorithms, providing
numerous examples of determining the computational complexity of various algorithms.
Finally, we discuss some of the most common time complexities that 
occur in the study of algorithms.


%%%%%%%%%%%%%%%%%%%%%%%%%%%%%%%%%%%%%%%%%%%%%%%%%%%%%%%%%%%%%%%%%%%%%
\section{Asymptotic Notation}\label{sec:asymptotic}
{\em Asymptotic notation}\index{asymptotic notation}
 is used to express and compare the growth rate of functions.
In our case, the functions will typically represent the running time of algorithms.
We will define the asymptotic notations in terms of nonnegative functions.
You will find more general definitions of these notations in other books,
but they are more complicated, more difficult to understand, and harder to work with.  
These added difficulties are a result of the possibility of the functions involved
being negative.
But the main reason for our use of the notations is to express the running time 
of algorithms.  Since the running time of an algorithm is always nonnegative, 
there is really no good reason to use the more cumbersome definitions.
We will focus on the notations most commonly used in the analysis of algorithms.

Asymptotic notation allows us to express the behavior of a function as the input
approaches infinity.  In other words, it is concerned about what happens to $f(n)$
as $n$ gets larger, and is not concerned about the value of $f(n)$ for small values of $n$.


We will define four of the most commonly used notations (an allude to the definition
of a fifth), providing a few brief examples of each.  We will then discuss some of
the most important and useful properties of these notations.  Finally, we will
present many more detailed examples.  

\subsection{The Notations}
We begin with the most commonly used of the notations: {\em Big-O ($O$)}. We will define it and give several examples
of its use. We will then present {\em Big-Omega ($\Omega$)}, {\em Big-Theta ($\Theta$)}, and {\rm little-o ($\omega$)} notations,
providing definitions and examples, and discussing how the notations relate to each other.

\newpage % Awkward break in the middle of the definition if I don't do this.

\begin{df}[Big-O]\index{Big-O}\index{a asymptotic@$O(f(n))$ (Big-O)}
Let $f$ be a nonnegative function. 

\medskip
\noindent
\begin{minipage}{.4\textwidth}
We say that $f(n)$ is {\bf Big-O} of $g(n)$, written as 
$f(n) = O(g(n))$,
iff there are positive constants $c$ and $n_{0}$ such that
$$f(n)\leq c\,g(n)\mbox{ for all }n \ge n_{0}.$$
If $f(n)=O(g(n))$, $f(n)$ grows no faster than $g(n)$. 
In other words, $g(n)$ is an {\bf asymptotic upper bound} (or just {\bf upper bound})
on $f(n)$.
\end{minipage}
\hfill
\begin{minipage}{.55\textwidth}
\sputfig{fig/c3-29.eps}{.55}
\end{minipage}
\end{df}

\begin{rem}
The ``='' in the statement ``$f(n)=O(g(n))$'' should be read and thought of as ``is'', not ``equals.''
You can think of it as a one-way equals.  So saying $f(n)=O(g(n))$ is {\em not}
the same thing as saying $O(g(n))=f(n)$, for instance 
(with the latter statement not really making sense).

An alternative notation is to write $f(n)\in O(g(n))$  instead of $f(n)=O(g(n))$. 
It turns out that $O(g(n))$ is actually the set of all functions that grow no 
faster than $g(n)$,
so the set notation is actually in some sense more correct.
The ``='' notation is used because it comes in handy when doing algebra.
You can essentially think of these as being two different notations ($=$ and $\in$) 
for the same thing.
Similar statements are true for the other asymptotic notations.
\end{rem}

\begin{exa}
Prove that $n^2+n=O(n^3)$.
\begin{sol}
Here, we have $f(n) = n^2+n$, and $g(n) = n^{3}$.
Notice that if $n\geq 1$, $n\leq n^3$ and $n^2\leq n^3$.
Therefore, 
$$n^2+n\leq n^3+n^3=2 n^3$$
Thus, 
$$n^2+n\leq 2 n^3 \mbox{ for all } n\geq 1.$$
Thus, we have shown that $n^2+n=O(n^3)$ by definition of Big-O, with $n_0=1$, and $c=2$.
\end{sol}
\end{exa}

The following fact is a generalization of what was used in the previous example.
It is used often in proofs involving asymptotic notation.
\begin{thm}
If $a$ and $b$ are real numbers with $a\leq b$, 
then $n^a\leq n^b$ whenever $n\geq 1$. 
\begin{pf}
We will not provide a proof, but it should be fairly clear intuitively that this is true.
If you cannot see why this is true, 
you should work out a few examples to convince yourself.
\end{pf}
\end{thm}

Sometimes the easiest way to prove that $f(n)=O(g(n))$ is to
take $c$ to be the sum of the {\em positive} coefficients of $f(n)$, 
although this trick doesn't always work.
We can usually easily eliminate the lower order terms with negative coefficients if we make 
the appropriate assumption.  Let's see how to do this in the next few examples.

\begin{exa}\label{exa:firstbigO}
Prove that $3n^{3}- 2n^2 + 13 n -15=O(n^3)$.
\begin{sol} 
First, notice that if $n\geq 0$, then $-2n^2-15\leq 0$, so 
$$3n^{3}- 2n^2 + 13 n -15\leq 3n^{3}+13n.$$
Next, if $n\geq 1$, then $13n\leq 13n^3$. 
Therefore if $n\geq 1$, $$3n^{3}+13n\leq 3n^3+13n^3=16n^3.$$
Also notice that if $n\geq 1$, then $n\geq 0$.
Thus, our first step is still valid if we assume $n\geq 1$ since $n\geq 1$
is a stronger condition than $n\geq 0$.  
Putting this all together, if we assume $n\geq 1$, then
\begin{eqnarray*}
3n^{3}- 2n^2 + 13 n -15 &\leq&3n^{3} +13 n \\
&\leq&3n^{3}+ 13 n^3 \\
&=&16n^3.
\end{eqnarray*} 
Since we have shown that $3n^{3}- 2n^2 + 13 n -15\leq 16n^3$ for all $n\geq 1$,
we have proven that $3n^{3}- 2n^2 + 13 n -15=O(n^3)$.  

We used $n_0=1$ and $c=16$ in our proof.  It is not necessary to explicitly point
this out in our proof, though.  We only do so to help you see the connection between
the proof and the definition of Big-O.
\end{sol}
\end{exa}


\begin{exa}\label{exa:bigoone}
Prove that $5n^{2}- 3n +20=O(n^2)$.
\begin{sol} 
If $n\geq 1$,
\begin{eqnarray}
\label{bigoone:step1}5n^{2} - 3n + 20 &\leq& 5n^{2} + 20 \\ 
\label{bigoone:step2}&\leq&5n^2+20n^2\\ 
&=&25n^2.
\end{eqnarray}
Since $5n^{2} - 3n + 20 \leq 25n^2$ for all $n\geq 1$,
$5n^{2} - 3n +20=O(n^2)$.

In this proof we used $c=25$ and $n_0=1$.
\end{sol}
\end{exa}

\begin{ques}
Answer the following questions related to Example~\ref{exa:bigoone}.
\begin{lenum}
\item
What allowed us to eliminate the $-3n$ term in step~\ref{bigoone:step1}?
\blanklinetwo
\item
What is the justification for step~\ref{bigoone:step2}?
\blanklinefill
\end{lenum}
\begin{answer}
(a) Since we assumed that $n\geq 1$, $-3n$ is certainly negative.  In other words,
$-3n\leq 0$.  That's why in the first step we could say that $5n^2-3n+20\leq 5n^2+20$.
(b) We used the fact that $20\leq 20n^2$ whenever $n\geq 1$.
If either of these solutions is not clear to you, 
you need to brush up on your algebra.
\end{answer}
\end{ques}


\begin{eval}
Prove that $4n^{2}- 12n +10=O(n^2)$.
\begin{ssol}
If $n\geq 1$, $4 n^{2} - 12n + 10 \leq 
4 n^{2} - 12n^2 + 10n^2 =2n^2.$
Therefore,  $4n^{2} - 12n +10=O(n^2)$.
\end{ssol}
\evallinetwo
\begin{answer}
This is incorrect.  
It is {\em not true} that $-12n\leq -12n^2$ when $n\geq 1$.
(If this isn't clear to you after thinking about it for a few minutes, you may need to
do some algebra review.)
In fact, that error led to the statement
$4 n^{2} - 12n + 10\leq 2n^2$ which cannot possibly be true as $n$ gets larger
since it would require that $2n^2-12n+10\leq 0$.
This is not true as $n$ gets larger.  
In fact, when $n=10$, for instance, it is clearly not true.
But it {\em is true} that $-12n<0$ when $n\geq 0$, so instead of replacing it
with $-12n^2$, it should be replaced with $0$ as in previous examples.
\end{answer}
\end{eval}

\begin{rem}
The values of the constants used in the proofs do not need to be the best possible. 
For instance, if you can show that $f(n)\leq 345\, g(n)$ for all $n\geq 712$, then $f(n)=O(g(n))$.
It doesn't matter whether or not it is actually true that 
$f(n)\leq 3\, g(n)$ for all $n\geq 5$.
\end{rem}

\begin{ques}
Answer each of the following questions related to Example~\ref{exa:bigoone}.  
Include a brief justification.
\begin{lenum}
\item
Could we have used $c=50$ in the proof?

\noindent
\answerlinetwo
\item
Could we have used $c=2$ in the proof?

\noindent
\answerlinetwo
\item
Could we have used $n_0=100$ in the proof?

\noindent
\answerlinetwo
\item
Could we have used $n_0=0$ in the proof?

\noindent
\answerlinetwo
\end{lenum}
\begin{answer}
(a) Sure.  Add the final step of $25n^2\leq 50n^2$ to the algebra in the proof.
In fact, any number above $25$ can easily be used.
Some values under 25 can also be used, but they would require a modification
of the algebra used in the proof.  
The bottom line is that there is generally no `right' value to
use for $c$.  If you find a value that works, then it's fine.
(b)  Clearly not.  For this to work, we would need
$5n^2-3n+20<2n^2$ to hold as $n$ increases towards $\infty$.
But this would imply that $3n^2-3n+20<0$.
But when $n\geq 1$, $3n^2$ is positive and larger than $3n$, so 
$3n^2-3n+20>0$.
(c) Sure.  The proof used the fact that the inequality is true when $n\geq 1$,
so it is clearly also true if $n\geq 100$.  And the definition of Big-O does
not require that we use the smallest possible value for $n_0$.
(d) No.  We would need a constant $c$ such that $5\cdot 0^2-3\cdot 0 + 20 = 20\leq 0= c\cdot 0^2$,
which is clearly impossible.
\end{answer}
\end{ques}

\newpage % another awkward one

\begin{exer}\label{exer:firstbigoh}
Prove that $5n^{5} -4n^4 + 3n^3 - 2n^2 +n=O(n^5)$.
(Hint: Use the same techniques you saw in Example~\ref{exa:bigoone}.)

\vspace*{3in}
\begin{answer}
If $n\geq 1$,
\begin{eqnarray*} 
5n^{5} -4n^4 + 3n^3 - 2n^2 +n
&\leq&5n^{5} + 3n^3 +n\\
&\leq&5n^{5} + 3n^5 +n^5\\
&=&9n^5.
\end{eqnarray*}
Therefore, $5n^{5} -4n^4 + 3n^3 - 2n^2 +n=O(n^5)$.
\end{answer}
\end{exer}

\begin{ques}
What values did you use for $n_0$ and $c$ in your solution to 
Exercise~\ref{exer:firstbigoh}?

\noindent
$n_0=\blankline{.75},\, \, \,  c=\blankline{.75}$
\begin{answer}
We used $n_0=1$ and $c=9$. 
Your values for $n_0$ and $c$ may differ.  
This is O.K. if you have the correct algebra to back it up.
\end{answer}
\end{ques}


Things are not always so easy.  How would you show that 
$(\sqrt{2})^{\log n}+\log^2 n + n^4=O(2^n)$?
Or that $n^2=O(n^2-13n+23)$?
In general, we simply (or in some cases with much effort) 
find values $c$ and $n_0$ that work. 
This gets easier with practice.

Big-O is a notation to express the idea that one function is an upper bound for 
another function.  The next notation allows us to express the opposite idea---that one
function is a {\em lower bound} for another function.

\begin{df}[Big-Omega]\index{Big-Omega}\index{a asymptotic@$\Omega(f(n))$ (Big-Omega)}
Let $f$ and $g$ be nonnegative functions.

\medskip
\noindent
\begin{minipage}{.4\textwidth}
We say that $f(n)$ is {\bf Big-Omega} of $g(n)$, written as 
$f(n) = \Omega(g(n))$,
iff there are positive constants $c$ and $n_{0}$ such that
$$c\,g(n) \leq f(n) \mbox{ for all } n \geq n_0.$$
When we say $f(n)=\Omega(g(n))$, it means that 
 $f(n)$ grows no slower than $g(n)$. In other words,
$g(n)$ is an {\bf asymptotic lower bound} (or just {\bf lower bound}) 
on $f(n)$.
\end{minipage}
\hfill
\begin{minipage}{.55\textwidth}
\sputfig{fig/c3-30.eps}{.55}
\end{minipage}
\end{df}


\begin{exa} 
Prove that $n^3+4\,n^2=\Omega(n^2)$.
\begin{pf}
Here, we have $f(n) = n^3+4\,n^2$, and $g(n) = n^2$.
It is not too hard to see that if $n\geq 1$, 
$$n^2\leq n^3\leq n^3+4\,n^2$$
Therefore, 
$$1\,n^2\leq n^3+4\,n^2 \mbox{ for all } n\geq 1$$
so $n^3+4\,n^2=\Omega(n^2)$ by definition of $\Omega$, with $n_0=1$, and $c=1$.
\end{pf}
\end{exa}


\begin{exer}\label{exer:omegaone}
Prove that $4n^2+n+1=\Omega(n^2)$. 
(This one should be really easy---follow the technique from the previous example and
don't over think it.)

\vspace*{1.3in}
\begin{answer}
Since $4n^2\leq 4n^2+n+1$ for $n\geq 0$, $4n^2=\Omega(n^2)$.
\end{answer}
\end{exer}

\begin{ques}
What values did you use for $n_0$ and $c$ in your solution to 
Exercise~\ref{exer:omegaone}?

\noindent
$n_0=\blankline{.75},\, \, \,  c=\blankline{.75}$
\begin{answer}
We used $c=4$ and $n_0=0$.  You might have used $n_0=1$ or some other positive value.
As long as you chose a positive value for $n_0$, it works just fine.  You could have also 
used any value for $c$ larger than $0$ and at most $4$.
\end{answer}
\end{ques}

Proving that $f(n)=\Omega(g(n))$ often requires more thought than proving that $f(n)=O(g(n))$.
Although the lower-order terms with positive coefficients can be easily dealt with,
those with negative coefficients make things a bit more complicated.
Often, we have to pick $c<1$.  
A good strategy is to pick a value of $c$ that you think will
work, and determine which value of $n_0$ is needed.
Being able to do some algebra helps.
As it turns out, we won't have to worry a whole lot about this, though.  We will
see a different technique to prove bounds shortly that, when it works, 
makes things much easier.

Our third notation allows us to express the idea that two functions grow at the 
same rate.

\begin{df}[Big-Theta]\index{Big-Theta}\index{a asymptotic@$\Theta(f(n))$ (Big-Theta)}
Let $f$ and $g$ be nonnegative functions.

\medskip
\noindent
\begin{minipage}{.45\textwidth}
We say that $f(n)$ is {\bf Big-Theta} of $g(n)$, written as 
$f(n) = \Theta(g(n))$,
iff there are positive constants $c_1$, $c_2$ and $n_{0}$ such that
$$c_1\,g(n)\leq f(n)\leq c_2\,g(n)\mbox{ for all }n \ge n_{0}.$$
If $f(n)=\Theta(g(n))$,
 $f(n)$ grows at the same rate as $g(n)$. In other words,
$g(n)$ is an {\bf asymptotically tight bound} (or just {\bf tight bound}) on $f(n)$.
\end{minipage}
\hfill
\begin{minipage}{.5\textwidth}
\sputfig{fig/c3-31.eps}{.45}
\end{minipage}
\end{df}

\begin{exa}
Prove that $n^2+5n+7=\Theta(n^2)$
\begin{pf}
When $n\geq 1$, $n^2+5\,n+7\leq n^2+5\,n^2+7n^2\leq 13n^2.$

\noindent
When $n\geq 0$,  $n^2\leq n^2+5\,n+7$.

\noindent
Combining these, we can see that when $n\geq 1$, 
$$n^2\leq n^2+5\,n+7\leq 13\,n^2,$$
so $n^2{+}5\,n{+}7=\Theta(n^2)$ by definition of $\Theta$, 
with $n_0{=}1$, $c_1{=}1$, and $c_2{=}13$.
\end{pf}
\end{exa}

\begin{ques}
In the previous example, we combined two inequalities.  One of them assumed $n\geq 0$, the other
assumed that $n\geq 1$.  In the combined inequality, we said it held if $n\geq 1$.  
Is that really O.K., or did we make a subtle error?

\noindent
\answerlinetwo
\begin{answer}
It is O.K.  Since the second inequality holds when $n\geq 0$, it also holds when $n\geq 1$.

In general, when you want to combine inequalities that contain two different assumptions, 
you simply make the more restrictive assumption.  In this case, $n\geq 1$ is more restrictive
than $n\geq 0$.  In general, if you have assumptions $n\geq a$ and $n\geq b$, then to combine
the results with these assumptions, you assume $n\geq max(a,b)$.
\end{answer}
\end{ques}


Using the definition of $\Theta$ can be inconvenient since it involves a double inequality.  Luckily, the
following theorem provides us with an easier approach. 

\begin{thm}\label{thm:theta}
If $f$ and $g$ are nonnegative functions, then $f(n) = \Theta(g(n))$ if and only if $f(n) = O(g(n))$ and $f(n) = \Omega(g(n))$.
\begin{pf}
The result follows almost immediately from the definitions.  We leave the details to the reader.
\end{pf}
\end{thm}

This theorem implies that no new strategies are necessary for $\Theta$ proofs 
since they can be split into two proofs---a Big-O proof and a $\Omega$ proof.
Let's see an example of this approach.

\begin{exa}\label{exa:secondtheta}
Show that $\frac{1}{2} n^{2} + 3\,n=\Theta(n^{2})$
\begin{pf}
Notice that if $n\geq 1$, 
$$\frac{1}{2} n^{2}+3\,n\leq \frac{1}{2}n^{2}+3\,n^{2}=\frac{7}{ 2}n^2,$$
so  ${\displaystyle \frac{1}{2}n^2 +3\,n=O(n^2)}$.
Also, when $n\geq 0$, 
$$ \frac{1}{2}n^2 \leq \frac{1}{2}n^2 +3n,$$
so ${\displaystyle \frac{1}{2}n^2 +3\,n=\Omega(n^2)}$.
Since 
$\frac{1}{2}n^2 +3\,n=O(n^2)$ and
$\frac{1}{2}n^2 +3\,n=\Omega(n^2)$, then by Theorem~\ref{thm:theta},
${\displaystyle \frac{1}{2}n^2 +3\,n=\Theta(n^2)}$
\end{pf}
\end{exa}

How do you use asymptotic notation to express that $f(n)$ grows slower than $g(n)$?
Saying $f(n)=O(g(n))$ doesn't work, because that only tells us that $f(n)$ grows 
{\em no faster than} $g(n)$.  
It might grow slower, but it also might grow at the same rate.
With the notation we have, the best way to express this idea is to say 
that $f(n)=O(g(n))$ and $f(n)\neq \Theta(g(n))$.  
But that is awkward.  Let's learn a new notation for this instead.
For technical reasons that we won't get into, this notation has to be defined
somewhat differently than the others.

\begin{df}\index{little-O}\index{a asymptotic@$o(f(n))$ (little-o)}
Let $f$ and $g$ be nonnegative functions, with $g$ being eventually non-zero.
We say that $f(n)$ is {\bf little-o} of $g(n)$, written $f(n)=o(g(n))$ iff
$$\lim_{n\rightarrow \infty}\dfrac{f(n)}{g(n)}=0.$$ 
If $f(n)=o(g(n))$, $f(n)$ grows {\bf asymptotically slower} than $g(n)$. 
\end{df}

\vspace*{-.1in}

\begin{exa}
You should be able to convince yourself that $3n+2=o(n^2)$, but $3n+2\neq o(n)$.
Similarly, $n^2+n + 4=o(n^3)$ and $n^2+n + 4=o(n^4)$, but 
$n^2+n + 4\neq o(n^2)$ and $n^2+n + 4\neq o(n)$.

If you are not comfortable with limits you can still convince yourself of these
statements by thinking of the informal definition.  
For instance, $n^2+n + 4$ grows slower than $n^3$ so $n^2+n + 4=o(n^3)$.
On the other hand, $n^2+n + 4$ grows at the same rate (so {\em not} slower than)
$n^2$, so $n^2+n + 4\neq o(n^2)$.
\end{exa}

\vspace*{-.1in}

\begin{ques}
Why do we require that $g(n)$ be eventually non-zero in the definition of little-o?

\noindent
\answerline
\begin{answer}
$g(n)$ appears in the denominator of a fraction.  If at some point it does not become 
(and remain) non-zero, the limit in the definition will be undefined.
If you never took a calculus course and are not that familiar with limits, 
do not worry a whole lot about this subtle point.
\end{answer}
\end{ques}

Little-omega ($\omega$) can be defined similarly to little-o, 
but the value of the limit is $\infty$ instead of $0$.
\index{little-omega}\index{little-omega}\index{a asymptotic@$\omega(f(n))$ (little-omega)}
We won't use $\omega$ very often.

\begin{ques}
Big-O notation is analogous to $\leq$ in certain ways. 
If so, what would be the similar analogies for $o$ and $\omega$?

\noindent
\answerline
\begin{answer}
$o$ is like $<$ and $\omega$ is like $>$.
\end{answer}
\end{ques}

\begin{rem}
\begin{itemize}
\item
It is important to remember that a $O$-bound is only an {\bf upper bound},
and that it may or may not be a tight bound.
So if $f(n)=O(n^{2})$, it is possible that $f(n)=3n^2+4$, $f(n)=\log n$, or
any other function that grows no faster than $n^2$.  But we also know that 
$f(n)\neq n^3$ or any other function that grows faster than $n^2$.
\item
Conversely, a $\Omega$-bound is only a {\bf lower bound}.  
Thus, if $f(n)=\Omega(n\log n)$, it might be the case that 
$f(n)=2^n$, but we know that $f(n)\neq 3n$, for instance.
\item
Unlike the others, $\Theta$-bounds are precise.  So, if
$f(n)=\Theta(n^2)$, then we know that $f$ has quadratic growth rate.
It might be that $f(n)=3n^2$,  $2n^2-43\,n-4$, or even $n^2+n\log n$.
But we are certain that the fastest growing term of $f$ is $c\, n^2$ for some constant $c$.
\end{itemize}
\end{rem}

\begin{ques}
Answer the following questions about the asymptotic notations.
\begin{lenum}
\item
If $f(n)=\Theta(g(n))$, is it possible that $f(n)=o(g(n))$?  Explain.

\noindent
\blanklinetwo
\item
If $f(n)=O(g(n))$, is it possible that $f(n)=o(g(n))$?  Explain.

\noindent
\blanklinetwo
\item
If $f(n)=O(g(n))$, is it {\em certain} that $f(n)=o(g(n))$?  Explain.

\noindent
\blanklinetwo
\item
If $f(n)=o(g(n))$, is it possible that $f(n)=O(g(n))$?  Explain.

\noindent
\blanklinetwo
%\item
%If $f(n)=\Omega(g(n))$, is it {\em certain} that $f(n)=\Theta(g(n))$?  Explain.\\
%\blanklinetwo
\end{lenum}
\begin{answer}
(a) No.  If $f(n)=\Theta(g(n))$, $f$ and $g$ grow at the same rate.  But 
$f(n)=o(g(n))$ expresses the idea that $f$ grows slower than $g$.
It is impossible for $f$ to grow at the same rate as $g$ and slower than $g$.
(b) Yes!  If $f$ grows no faster than $g$, then it is possible that it grows slower.
For instance, $n=O(n^2)$ and $n=o(n^2)$ are both true.
(c) No.  If $f$ and $g$ grow at the same rate, then $f(n)=O(g(n))$, but 
$f(n)\neq o(g(n))$. For instance, $3n=O(n)$, but $3n\neq o(n)$. 
(d) Yes.  In fact, it is guaranteed!  If $f$ grows slower than $g$, then
$f$ grows no faster than $g$. 
%(e) No.  It is possible, but not certain.  For instance, 
%$n^2=\Omega(n)$, but $n^2\neq\Theta(n)$.
\end{answer}
\end{ques}



\begin{eval}
Let $a_0,\ldots, a_k\in \reals$, where $a_k>0$.
Prove that $a_kn^k+a_{k-1}n^{k-1}+\cdots+a_1n+a_0=O(n^k)$.
\begin{ssolenum}
\item
We can first eliminate all of the constants since they become irrelevant as $n$ grows large enough.
This leaves us with $n^k+n^{k-1}+\cdots+n=O(n^k)$.  Next we can eliminate
all terms growing slower than $n^k$, since they also become irrelevant as $n$ grows.
This leaves us with $n^k=O(n^k)$, and since they are the same, they are effectively theta
of each other, and by definition, anything that is theta of something is also omega and $O$,
so we can correctly say that $n^k=O(n^k)$, thus proving that 
$a_kn^k+a_{k-1}n^{k-1}+\cdots+a_1n+a_0=O(n^k)$.

\noindent
\evallinethree
\item
Let $c=\dsum_{i=0}^k |a_i|$.  
Then if $n\geq 1$, 
\begin{eqnarray*}
a_kn^k+a_{k-1}n^{k-1}+\cdots+a_1n+a_0&\leq&|a_k|n^k+|a_{k-1}|n^{k-1}+\cdots+|a_1|n+|a_0|\\
&\leq&|a_k|n^k+|a_{k-1}|n^k+\cdots+|a_1|n^k+|a_0|n^k\\
&\leq&\dsum_{i=0}^k |a_i|n^k=cn^k.
\end{eqnarray*}
Therefore, $a_kn^k+a_{k-1}n^{k-1}+\cdots+a_1n+a_0=O(n^k)$.

\noindent
\evallinethree
\end{ssolenum}
\begin{answer}
\begin{solevalenum}
\item
Although this proof sounds somewhat reasonable, it is way too informal and convoluted.  Here are
some of the problems.
\begin{enumerate}
\item
This student misunderstands the concept behind `ignoring the constants.'
We can ignore the constants {\em after we know that $f(n)=O(g(n))$}.  We can't ignore them
in order to prove it.
\item
The phrase `become irrelevant' (used twice) is not precise.  
We have developed mathematical notation for
a reason---it allows us to make statements like these precise.
It's kind of like saying that a certain car costs `a lot'.  What is `a lot'?
Although \$30,000 might be a lot for most of us, people with a lot more money
than I have might not think that \$500,000 is a lot.
\item
The phrase `This leaves us with $n^k+n^{k-1}+\cdots+n=O(n^k)$' is odd.  What precisely do they mean?
That this is true or that this is what we need to prove now?  In either case, it is incorrect.
Similarly for the second time they use the phrase `This leaves us with'.
\item
The second half of the proof is unnecessarily convoluted.  They essentially are claiming that their
proof has boiled down to showing that $n^k=O(n^k)$.  
To prove this, they use an incredibly drawn out, yet vague, explanation that is 
in a single unnecessarily long sentence.  Why are they even bringing $\Theta$ and $\Omega$ into this proof?  
Why don't they just say something like `since $n^k\leq 1n^k$ for all $n\geq 1$, $n^k=O(n^k)$'?  
I believe the answer is obvious:  they don't really understand what they are doing here.  
They clearly have a vague understanding of the notation, but they don't understand the 
formal definition.
\end{enumerate}
The bottom line is that this student understands that the statement they needed to prove is correct,
and they have a vague sense of {\em why} it is true, but they did not have a clear understanding
of how to use the definition of Big-O to prove it.
The most important thing to take away from this example is this:  
{\em Be precise, use the notation and
definitions you have learned, and if your proofs look a lot different than those in the book, you
might question whether or not you are on the right track.}
\item
This proof is correct.
\end{solevalenum}
\end{answer}
\end{eval}

\begin{exer}
Assume that $f(n)=O(n^2)$ and $g(n)=O(n^3)$.
What can you say about the relative growth rates of $f(n)$ and $g(n)$?
In particular, does $g(n)$ grow faster than $f(n)$?

\noindent
\answerlinethree
\begin{answer}
We cannot say anything about the relative growth rates of $f(n)$ and $g(n)$
because we are only given upper bounds for each. 
It is possible that $f(n)=n^2$ and $g(n)=n$, so that 
$f(n)$ grows faster, or vice-versa.  They could also both be $n$.
\end{answer}
\end{exer}

Keep in mind that asymptotic notation only allows you
to compare the asymptotic behavior of functions.  
Except for $\Theta$-notation, it only provides
a bound on the growth rate.  For instance, knowing that $f(n)=O(g(n))$ only tells
you that $f(n)$ grows no faster than $g(n)$.  
It is possible that $f(n)$ grows a lot slower than $g(n)$.



\begin{exer}
Let's test your understanding of the material so far. Answer each of the following
true/false questions, giving a very brief justification/counterexample.
Justifications can appeal to a definition and/or theorem.
For counterexamples, use simple functions. For instance, $f(n)=n$ and $g(n)=n^2$.
\begin{lenum}
\itemtf
If $f(n)=O(g(n))$, then $f(n)$ grows faster than $g(n)$
\bigbreak
\itemtf
If $f(n)=\Theta(g(n))$, then $f(n)$ grows faster than $g(n)$ 
\bigbreak
\itemtf
If $f(n)=O(g(n))$, then $f(n)$ grows at the same rate as $g(n)$
\bigbreak
\itemtf
If $f(n)=\Omega (g(n))$, then $f(n)$ grows faster than $g(n)$ 
\bigbreak
\itemtf
If $f(n)=O(g(n))$, then $f(n)=\Omega(g(n))$
\bigbreak
\itemtf
If $f(n)=\Theta(g(n))$,  then $f(n)=O(g(n))$
\bigbreak
\itemtf
If $f(n)=O(g(n))$,  then $f(n)=\Theta(g(n))$
\bigbreak
\itemtf
If $f(n)=O(g(n))$,  then $g(n) = O(f(n))$
\end{lenum}
\vspace*{.25in}
\begin{answer}
(a) F. This is saying that $f(n)$ grows {\em no} faster than $g(n)$.
(b) F. They grow at the same rate.
(c) F. $f(n)$ might grow slower than $g(n)$.  For instance, $f(n)=n$ and $g(n)=n^2$.
(d) F. They might grow at the same rate.  For instance, $f(n)=g(n)=n$.
(e) F. If $f(n)=n$ and $g(n)=n^2$, $f(n)=O(g(n))$, but $f(n)\neq \Omega(g(n))$.
(f) T. By Theorem~\ref{thm:theta}.
(g) F. If $f(n)=n$ and $g(n)=n^2$, $f(n)=O(g(n))$, but $f(n)\neq \Theta(g(n))$. 
(h) F. If $f(n)=n$ and $g(n)=n^2$, $f(n)=O(g(n))$, but $g(n)\neq O(f(n))$.
\end{answer}
\end{exer}

\subsection{Properties of the Notations}
There are a lot of properties that hold for Big-O, $\Theta$ and $\Omega$ notation
(and $o$ and $\omega$ as well, but we won't focus on those ones in this section).
We will only present a few of the most important ones.  We provide proofs for
some of the results.  
The rest can be proven without too much difficulty using the definitions of the notations.

Before we present the properties, it might be useful to think about the properties of
things you are already familiar with.  
For instance, given real numbers $x$, $y$ and $z$, 
you know that if $x\leq y$ and $y\leq z$, then $x\leq z$.  This is just the
transitive property of $\leq$.  
Similarly, you know that if $x\leq y$, then $a\,x\leq a\,y$ for any positive constant $a$.
You can think of Big-O notation as being like $\leq$, $\Theta$ notation as being like $=$,
and $\Omega$ notation as being like $\geq$.  Many of the properties of $\leq$, $=$ and $\geq$
that you are already familiar with have an analog with Big-O, $\Theta$, and $\Omega$ notation.
But you need to be careful because the analogies are not exact.  For instance, constants cannot
be ignored with inequalities but can be ignored when using asymptotic notation.

\begin{thm}\label{thm:asyprops}
The transitive property holds for Big-O, $\Theta$, and $\Omega$.  That is,
\begin{itemize}
\item If $f(n)=O(g(n))$ and $g(n) =O(h(n)),$ then  $f(n) = O(h(n))$
\item If $f(n)=\Theta(g(n))$ and $g(n) = \Theta(h(n))$, then $f(n) = \Theta(h(n))$
\item If $f(n)=\Omega(g(n))$ and $g(n) = \Omega(h(n))$, then $f(n) = \Omega(h(n))$
\end{itemize}
\begin{pf}
You will prove the transitive property of Big-O in Exercise~\ref{exer:bigoTranProof}.
The proofs of the other two are very similar.
\end{pf}
\end{thm}

Theorem~\ref{thm:asyprops} is pretty intuitive.  
For instance,
when applied to Big-O notation, Theorem~\ref{thm:asyprops} is
essentially stating that if $g(n)$ is an upper bound on $f(n)$ and
$h(n)$ is an upper bound on $g(n)$, then $h(n)$ is an upper bound for $f(n)$.
Put another way, if $f(n)$ grows no faster than $g(n)$ and $g(n)$ grows no faster
than $h(n)$, then $f(n)$ grows no faster than $h(n)$.
This makes perfect sense if you think about it for a few minutes. 

  
\begin{exa}
Let's take it for granted that $4n^2+3n+17=O(n^3)$ and $n^3=O(n^4)$
(both of which you should be able to easily prove at this point).
According to Theorem~\ref{thm:asyprops}, we can conclude that $4n^2+3n+17=O(n^4)$.
\end{exa}

\begin{thm}\label{thm:asyprops2}
Scaling by a constant factor

\noindent
If $f(n)= O(g(n))$, then for any $k > 0$, $k f(n) = O(g(n))$.
Similarly for $\Theta$ and $\Omega$.
\begin{pf}
We will give the proof for Big-O notation.  
The other two proofs are similar.
Assume $f(n)= O(g(n))$.  Then by the definition of Big-O, there are positive constants $c$ and
$n_0$ such that $f(n)\leq c\,g(n)$ for all $n\geq n_0$.
Thus, if $n\geq n_0$, 
$$k\,f(n)\leq k\,c\,g(n)=c^{\prime}g(n),$$
where $c^{\prime}=k\,c$ is a positive constant.
By the definition of Big-O, $k f(n) = O(g(n))$.
\end{pf}
\end{thm}

\begin{exa}
Example~\ref{exa:secondtheta} showed that $\frac{1}{2} n^{2} + 3\,n=\Theta(n^{2})$.
We can use Theorem~\ref{thm:asyprops2} to conclude that 
$n^{2} + 6\,n=\Theta(n^{2})$ since $n^{2} + 6\,n=2\left(\frac{1}{2} n^{2} + 3\,n\right).$
\end{exa}

Perhaps now is a good time to point out a related issue.  Typically, we do not
include constants inside asymptotic notations.  For instance, although it is technically
correct to say that $34n^3+2n^2-45n+5=O(5n^3)$ (or $O(50n^3)$, or any other constant
you care to place there), it is best to just say it is $O(n^3)$.  
In particular, $\Theta(1)$ may be preferable to $\Theta(k)$.

\begin{thm}\label{thm:asyprops3}
Sums\\
 If $f_{1}(n) = O(g_{1}(n))$ and $f_{2}(n) = O(g_{2}(n))$,
then 
$$f_{1}(n)+f_{2}(n) = O(g_{1}(n)+g_{2}(n))=O(max\{g_{1}(n),g_{2}(n)\}).$$
Similarly for $\Theta$ and $\Omega$.
\begin{pf}
We will prove the assertion for Big-O.
Assume $f_{1}(n) = O(g_{1}(n))$ and $f_{2}(n) = O(g_{2}(n))$.
Then there exists positive constants $c_1$ and $n_1$ such that for all $n\geq n_1$, 
$$f_1(n)\leq c_1g_1(n),$$
and there exists positive constants $c_2$ and $n_2$ such that for all $n\geq n_2$, 
$$f_2(n)\leq c_2g_2(n).$$
Let $c_0=\max\{c_1,c_2\}$ and $n_0=\max\{n_1,n_2\}$.
Since $n_0$ is at least as large as $n_1$ and $n_2$, then for all $n\geq n_0$, 
$f_1(n)\leq c_1g_1(n)$ and $f_2(n)\leq c_2g_2(n)$.
(If you don't see why this is, think about it.  This is a subtle but important step.)
Similarly, if $f_1(n)\leq c_1g_1(n)$, then clearly $f_1(n)\leq c_0g_1(n)$ since
$c_0$ is at least as big as $c_1$ (and similarly for $f_2$).
Then for all $n\geq n_0$, we have
\begin{eqnarray*}
f_1(n)+f_2(n)&\leq& c_1g_1(n)+c_2g_2(n)\\
&\leq& c_0g_1(n)+c_0g_2(n)\\
&\leq& c_0[g_1(n)+g_2(n)]\\
&\leq& c_0[\max\{g_1(n),g_2(n)\}+\max\{g_1(n),g_2(n)\}]\\
&\leq& 2c_0\max\{g_1(n),g_2(n)\}\\
&\leq& c\max\{g_1(n),g_2(n)\},
\end{eqnarray*}
where $c=2c_0$.
By the definition of Big-O, we have shown that $f_{1}(n)+f_{2}(n)=O(max\{g_{1}(n),g_{2}(n)\})$.
\end{pf}
Notice that in this proof we used $c=2\max\{c_1,c_2\}$ and $n_0=\max\{n_1,n_2\}$.
\end{thm}

Without getting too technical,
the previous theorem implies that you can upper bound
the sum of two or more functions by finding the upper bound of the fastest growing 
of the functions.  
Another way of thinking about it is if you ever have two or more functions inside Big-O notation,
you can simplify the notation by omitting the slower growing function(s).
It should be pointed out that there is a subtle point in this result about how to 
precisely define the maximum of two functions.  Most of the time the intuitive
definition is sufficient so we won't belabor the point.  

\begin{exa}
Since we have previously shown that $5n^2-3n+20=O(n^2)$ and that $3n^3-2n^2+13n-15=O(n^3)$,
we know that $(5n^2-3n+20) +(3n^3-2n^2+13n-15)=O(n^2+n^3)=O(n^3)$.
\end{exa}

\begin{thm}\label{thm:asyprops3b}
Products\\
 If $f_{1}(n) = O(g_{1}(n))$ and $f_{2}(n) = O(g_{2}(n))$,
then 
$$f_{1}(n)f_{2}(n) = O(g_{1}(n)g_{2}(n)).$$
Similarly for $\Theta$ and $\Omega$.
\end{thm}

\begin{exa}
Since we have previously shown that $5n^2-3n+20=O(n^2)$ and that $3n^3-2n^2+13n-15=O(n^3)$,
we know that $(5n^2-3n+20)(3n^3-2n^2+13n-15)=O(n^2n^3)=O(n^5)$.  
Notice that we could arrive at this same conclusion by multiplying the two polynomials
and taking the highest term.  However, this would require a lot more work than is necessary.
\end{exa}


The next theorem essentially says that if $g(n)$ is an upper bound on $f(n)$, then
$f(n)$ is a lower bound on $g(n)$.  This makes perfect sense if you think about it.

\begin{thm}\label{thm:asyprops4}
Symmetry (sort of)\\
 $f(n)=O(g(n))$ iff $g(n)=\Omega(f(n))$.
\end{thm}


It turns out that $\Theta$ defines an equivalence relation on the set of functions
from ${\bf Z}^{+}$ to ${\bf Z}^{+}$. That is, it defines a partition on these functions,
with two functions being in the same partition (or the same equivalence class)
if and only if they have the same growth rate.
But don't take our word for it.  You will help to prove this fact next.


\begin{fib}
Let $R$ be the relation on the set of functions from ${\bf Z}^{+}$
to ${\bf Z}^{+}$ such that $(f,g)\in R$ if and only if $f=\Theta(g)$.
Show that $R$ is an equivalence relation.
\begin{pf}
We need to show that $R$ is reflexive, symmetric, and transitive.

\noindent
{\bf Reflexive:}
Since $1\cdot f(n)\leq f(n) \leq 1\cdot f(n)$ for all $n\geq 1$, $f(n)=\Theta(f(n))$, so
$R$ is reflexive.

\noindent
{\bf Symmetric:} 
If $f(n)=\Theta(g(n))$, then there exist positive constants $c_1$, $c_2$, and $n_0$ such that
\blanklinefill

\noindent
This implies that
$$g(n)\leq \frac{1}{ c_1} f(n) \mbox{ and } g(n)\geq \frac{1}{ c_2} f(n) \mbox{ for all } n\geq n_0$$
which is equivalent to 
$$\blankline{1.8}\leq g(n) \leq\blankline{1.8}\mbox{ for all } n\geq n_0.$$
Thus $g(n)=\Theta(f(n))$, and $R$ is symmetric.

\noindent
{\bf Transitive:}
If $f(n)=\Theta(g(n))$, then there exist positive constants $c_1$, $c_2$, and $n_0$ such that
$$c_1 g(n)\leq f(n) \leq c_2 g(n) \mbox{ for all } n\geq n_0.$$
Similarly
if $g(n)=\Theta(h(n))$, then there exist positive constants $c_3$, $c_4$, and $n_1$ such that
\blanklinefill

\noindent
Then 
$$f(n)\geq c_1\, g(n) \geq c_1c_3\, h(n) \mbox{ for all } n\geq \max\{n_0,n_1\},$$
and 
$$f(n)\leq \blankline{.75}\, g(n) \leq \blankline{.75}\, h(n) \mbox{ for all } n\geq \blankline{.75} $$
$$\mbox{Thus, \, } \blankline{1} \leq f(n) \leq \blankline{1}  \mbox{ for all } n\geq \max\{n_0,n_1\}.$$
Since $c_1 c_3$ and $c_2 c_4$ are both positive constants, 
$f(n)=\blankline{1}$ by the definition of \blankline{.75}, so $R$ is \blankline{1}.
\end{pf}
\begin{answer}
$c_1 g(n)\leq f(n) \leq c_2 g(n) \mbox{ for all } n\geq n_0$;
$\frac{1}{ c_2} f(n)$;
$\frac{1 }{ c_1} f(n)$;
$c_3 h(n)\leq g(n) \leq c_4 h(n) \mbox{ for all } n\geq n_1.$;
$c_2$; $c_2 c_4$;
$\max\{n_0,n_1\}$;
$c_1 c_3 h(n)$;
$c_2 c_4 h(n)$;
$\Theta(h(n))$;
$\Theta$;
transitive;
\end{answer}
\end{fib}

\begin{exa}
The functions $n^2$, $3n^2-4n+4$,
$n^2+\log n$, and $3n^2+n+1$ are all $\Theta(n^2)$.  That is, they all have the 
same rate of growth and all belong to the same equivalence class. 
\end{exa}

\begin{exer}
Let's test your understanding of the material so far. Answer each of the following
true/false questions, giving a very brief justification/counterexample.
Justifications can appeal to a definition and/or theorem.
For counterexamples, use simple functions. For instance, $f(n)=n$ and $g(n)=n^2$.
\begin{lenum}
\itemtf
If $f(n)=O(g(n))$,  then $g(n) = \Omega(f(n))$ 
\bigbreak
\itemtf
If $f(n)=\Theta(g(n))$, then $f(n)=\Omega(g(n))$ and $f(n)=O(g(n))$
\bigbreak
\itemtf
If $f_1(n) = O(g_1(n))$ and $f_2(n) = O(g_2(n))$, 
then $f_1(n) + f_2(n) = O(max( g_1(n) , g_2(n)) )$   
\bigbreak
\itemtf
$f(n)=O(g(n))$  iff $f(n)=\Theta(g(n))$
\bigbreak
\itemtf
$f(n)=O(g(n))$  iff $g(n) = O(f(n))$
\bigbreak
\itemtf
$f(n)=O(g(n))$  iff $g(n) = \Omega(f(n))$
\bigbreak
\itemtf
$f(n)=\Theta(g(n))$ iff $f(n)=\Omega(g(n))$ and $f(n)=O(g(n))$
\bigbreak
\itemtf
If $f(n)=O(g(n))$ and $g(n) = O(h(n))$, then $f(n)=O(h(n))$
\end{lenum}
\vspace*{.25in}
\begin{answer}
(a) T. By Theorem~\ref{thm:asyprops4}.
(b) T. By Theorem~\ref{thm:theta}.
(c) T. By Theorem~\ref{thm:asyprops3}.
(d) F. The backwards implication is true, but the forward one is not.  For instance, 
if $f(n)=n$ and $g(n)=n^2$, clearly $f(n)=O(g(n))$, but $f(n)\neq \Theta(g(n))$.
(e) F. Neither direction is true.  For instance, 
if $f(n)=n$ and $g(n)=n^2$, $f(n)=O(g(n))$, but $g(n)\neq O(f(n))$.
(f) T. By Theorem~\ref{thm:asyprops4}.
(g) T. By Theorem~\ref{thm:theta}.
(h) T. By Theorem~\ref{thm:asyprops}.
\end{answer}
\end{exer}



%%%%%%%%%%%%%%%%%%%%%%%%%%%
\subsection{Proofs using the definitions}\label{subsec:alProof}

In this section we provide more examples and exercises that use the definitions to 
prove bounds.

The first example is annotated with comments (given in footnotes)
about the techniques that are used in many of these proofs.
We use the following terminology in our explanation.
By {\em lower order term} we mean a term that grows slower, and {\em higher order} means a term that 
grows faster.  The {\em dominating term} is the term that grows the fastest.  For instance, 
in $x^3+7\, x^2-4$, the $x^2$ term is a lower order term than $x^3$, and $x^3$ is the dominating term.
We will discuss common growth rates, including how they relate to each other, 
in Section~\ref{sec:growrates}.
%\ref{section:growthRates}
But for now we assume you know that
$x^5$ grows faster than $x^3$, for instance.


\begin{exa}\label{exa:boundExample}
Find a tight bound on $f(n)=n^8+7n^7-10n^5-2n^4+3n^2-17$.
\begin{sol}
We will prove that $f(n)=\Theta(n^8)$.
First, we will prove an upper bound for $f(n)$.  It is clear that 
when $n\geq 1$, 
\begin{eqnarray*}
n^8+7n^7-10n^5-2n^4+3n^2-17&\leq& n^8+7n^7+3n^2\, \, \, \, 
\footnote{We can upper bound any function by removing the lower order
terms with negative coefficients, as long as $n\geq 0$.}\\
&\leq& n^8+7n^8+3n^8\, \, \, \, 
\footnote{We 
can upper bound any function by replacing lower order terms 
that have
positive coefficients by the dominating term with the same coefficients.
Here, we must make sure that the dominating term is larger than the 
given term for all values of $n$ larger than some threshold $n_0$, and
we must make note of the threshold value $n_0$.
}\\
&=&11n^8
\end{eqnarray*}
Thus, we have 
$$f(n)=n^8+7n^7-10n^5-2n^4+3n^2-17\leq 11n^8 \mbox{ for all } n\geq 1,$$
and we have proved that $f(n)=O(n^8)$.

Now, we will prove the lower bound for $f(n)$.
When $n\geq 1$, 
\begin{eqnarray*}
n^8+7n^7-10n^5-2n^4+3n^2-17 &\geq& n^8-10n^5-2n^4-17\, \, \, \, 
\footnote{We can lower bound any function by removing the lower order
terms with positive coefficients, as long as $n\geq 0$.}\\ 
&\geq& n^8-10n^7-2n^7-17n^7\, \, \, \, 
\footnote{We can lower bound any function by replacing lower order terms with
negative coefficients by a sub-dominating term with the same coefficients.
(By sub-dominating, I mean one which dominates all but the dominating
term.)
Here, we must make sure that the sub-dominating term is larger than the 
given term for all values of $n$ larger than some threshold $n_0$, and
we must make note of the threshold value $n_0$.
Making a wise choice for which sub-dominating term to use is crucial
in finishing the proof.}\\
&=&n^8-29n^7
\end{eqnarray*}
Next, we need to find a value $c>0$ such that $n^8-29n^7\geq c n^8$.
Doing a little algebra, we see that this is equivalent to
$(1-c)n^8 \geq 29n^7$.  When $n\geq 1$, we can divide by $n^7$ and
obtain $(1-c)n\geq 29$. Solving for $c$ we obtain 
$$c\leq 1-\frac{29}{ n}.$$
If $n\geq 58$, then $c=1/2$ suffices.
We have just shown that if $n\geq 58$, then 
$$f(n)=n^8+7n^7-10n^5-2n^4+3n^2-17\geq \frac{1}{2}n^8.$$
Thus, $f(n)=\Omega(n^8)$.
Since we have shown that $f(n)=\Omega(n^8)$ and that 
$f(n)=O(n^8)$, we have shown that $f(n)=\Theta(n^8)$.
\end{sol}
\end{exa}

Let's see another example of a $\Omega$ proof.  You should note
the similarities between this and the second half of the proof in
the previous example.
\begin{exa}
Show that $(n \log n - 2\,n + 13)=\Omega(n \log n)$
\begin{pf}
We need to show that there exist positive constants $c$ and $n_0$
such that 
$$c\,n\log n \leq n\log n -2\,n + 13\mbox{ for all } n\geq n_0.$$
Since $n\log n -2\,n \leq n\log n -2\,n+13,$
we will instead show that 
$$c\,n\log n \leq n\log n -2\,n,$$
which is equivalent to
$$c\leq 1-\frac{2}{\log n}, \mbox{ when } n > 1.$$ 
If $n\geq 8$, then $2/(\log n)\leq 2/3$, and picking
$c=1/3$ suffices.
In other words, we have just shown that if $n\geq 8$, 
$$\frac{1}{3}\,n\log n \leq n\log n -2\,n.$$
Thus if $c=1/3$ and $n_0=8$, then for all $n\geq n_0$, we have
$$c\,n\log n \leq n\log n -2\,n \leq n\log n -2\,n + 13.$$
Thus $(n \log n - 2\,n + 13)=\Omega(n \log n)$.
\end{pf}
\end{exa}

\begin{fib}
Show that $\frac{1}{2} n^{2} - 3n = \Theta(n^{2})$
\begin{pf}
We need to find positive constants $c_1$, $c_2$, and $n_0$ such that
$$\blankline{1} \leq \frac{1}{2}n^2 - 3\,n \leq \blankline{1}\mbox{ for all } n\geq n_0$$
Dividing by $n^{2}$, we get  $c_{1} \leq \blankline{1} \leq c_{2}.$

Notice that if $n\geq 10$,  $\frac{1}{2} - \frac{3}{n} \geq \frac{1}{2} - \frac{3}{10} = \blankline{1.5},$
so we can choose $c_1 =1/5$. 
If $n\geq 10$, we also have that $\frac{1}{2} - \frac{3}{n}\leq \frac{1}{2},$
so we can choose $c_2=1/2$.
Thus, we have shown that 
$$\blankline{1} \leq \frac{1}{2}n^2 - 3\,n \leq \blankline{1} 
\mbox{ for all } n\geq \blankline{.5} .$$
Therefore, $\frac{1}{2} n^{2} - 3n =\Theta(n^{2})$.
\end{pf}
\begin{answer}
$c_{1}n^2$; 
$c_{2}n^2$; 
$\frac{1}{2} - \frac{3}{n}$;
$\frac{10-6}{20}=\frac{1}{5}$; 
$\frac{1}{5}n^2$; 
$\frac{1}{2}n^2$; 
 $10$.
\end{answer}
\end{fib}

\begin{ques}
In the previous proof, we claimed that if $n\geq 10$, 
$$\frac{1}{2} - \frac{3}{n} \geq \frac{1}{2} - \frac{3}{10}.$$
Why is this true?

\noindent
\answerlinethree
\begin{answer}
There are a few ways to think about this.  First, the larger $n$ is, the smaller $\frac{3}{n}$ is,
so a smaller amount is being subtracted.  But that's perhaps too fuzzy.  Let's look at it this way:
$$n\geq 10\Rightarrow 
 \frac{10}{3} \leq \frac{n}{3}
\Rightarrow \frac{3}{n}\leq \frac{3}{10}
\Rightarrow -\frac{3}{n}\geq -\frac{3}{10}
\Rightarrow \frac{1}{2}-\frac{3}{n}\geq \frac{1}{2}-\frac{3}{10}.
$$
\end{answer}
\end{ques}


\begin{exa}
Show that $(\sqrt{2})^{\log n}=O(\sqrt{n})$, where the base of the $\log$ is 2.
\begin{pf}
It is not too hard to see that 
$$(\sqrt{2})^{\log n}=n^{\log \sqrt{2}}=n^{\log 2^{1/2}}=
n^{\frac{1}{2}\log 2}=n^{\frac{1}{2}}=\sqrt{n}.$$
Thus it is clear that $(\sqrt{2})^{\log n}=O(\sqrt{n})$.
\end{pf}
\end{exa}

\begin{rem}
You may be confused by the previous proof.  It seems that we never showed that 
$(\sqrt{2})^{\log n}\leq c\sqrt{n}$ for some constant $c$.  
But we essentially did by showing that $(\sqrt{2})^{\log n}=\sqrt{n}$ since
this implies that 
$(\sqrt{2})^{\log n}\leq 1\sqrt{n}$. 
 
We actually proved something stronger than was required.
That is, since we proved the two functions are equal, it is in fact true that 
$(\sqrt{2})^{\log n}=\Theta(\sqrt{n})$.  But we were only asked to prove 
that $(\sqrt{2})^{\log n}=O(\sqrt{n})$.

In general, if you need to prove a Big-O bound, you may instead prove a $\Theta$ 
bound, and the Big-O bound essentially comes along for the ride.
\end{rem}

\begin{ques}
In our previous note we mentioned that if you prove a $\Theta$ bound, 
you get the Big-O bound for free.  
\begin{lenum}
\item
What theorem implies this?

\noindent
\answerline
\item
If we prove $f(n)=O(g(n))$, does that imply that $f(n)=\Theta(g(n))$?
In other words, does it work the other way around? Explain, giving an appropriate example.

\noindent
\answerlinetwo
\end{lenum}
\begin{answer}
(a) Theorem~\ref{thm:theta}.
(b) Absolutely not!  Theorem~\ref{thm:theta} requires that we also prove 
$f(n)=\Omega(g(n))$.  Here is a counterexample: $n=O(n^2)$, but $n\neq\Theta(n^2)$. 
So $f(n)=O(g(n))$ does not imply that $f(n)=\Theta(g(n))$.
\end{answer}
\end{ques}

\begin{exer}
Show that $n!=O(n^n)$. 
(Don't give up too easily on this one---the proof is very short and only uses
elementary algebra.)

\vspace*{2.25in}
\begin{answer}
Notice that when $n\geq 1$, 
$n!=1\cdot 2\cdot 3\cdots n \leq n\cdot n\cdots n = n^n.$
Therefore $n!=O(n^n)$ (We used $n_0=1$, and $c=1$.)
\end{answer}
\end{exer}


\begin{exa}
Show that $\log (n!) = O(n\log n)$
\begin{pf}
It should be clear that if $n\geq 1$, $n!\leq n^n$ (especially after 
completing the previous exercise). 
Taking $\log$s of both sides of that inequality,
we obtain 
$$\log n!\leq \log (n^n) = n\log n.$$
Therefore $\log n! = O(n\log n)$.
\end{pf}
The last step used the fact that $\log (f(n)^a)=a\log (f(n))$, a fact that we assume you 
have seen previously (but may have forgotten).
\end{exa}


Proving properties of the asymptotic notations is actually no more difficult than the rest of
the proofs we have seen.  You have already seen a few and helped write one.  
Here we provide one more example and then ask you to prove another result on your own.

\begin{exa}
Prove that if $f(n)=O(g(n))$ and $g(n)=O(f(n))$, then
$f(n)=\Theta(g(n))$.
\begin{pf}
If $f(n)=O(g(n))$, 
then there are positive constants $c_2$ and $n_{0}^{\prime}$ such that 
$$f(n)\leq c_2\,g(n)\mbox{ for all }n \ge n_{0}^{\prime}$$
Similarly, if $g(x)=O(f(x))$, 
then there are positive constants $c_1^{\prime}$ and $n_{0}^{\prime\prime}$
such that 
$$g(n)\leq c_1^{\prime}\,f(n)\mbox{ for all }n 
\ge n_{0}^{\prime\prime}.$$
We can divide this by $c_1^{\prime}$ to obtain 
$$\frac{1}{c_1^{\prime}}g(n)\leq f(n)\mbox{ for all }n 
\ge n_{0}^{\prime\prime}.$$
Setting $c_1=1/c_1^{\prime}$ and 
$n_0=\max\{n_0^{\prime},n_0^{\prime\prime}\}$, 
we have 
$$c_1 g(n)\leq f(n)\leq c_2\,g(n)\mbox{ for all }n \ge n_{0}.$$
Thus, $f(x)=\Theta(g(x)).$
\end{pf}
\end{exa}

\begin{exer}\label{exer:bigoTranProof}
Let $f(x)=O(g(x))$ and $g(x)=O(h(x))$. Show that $f(x)=O(h(x))$.
That is, prove Theorem~\ref{thm:asyprops} for Big-O notation.

\vspace*{3.5in}
\begin{answer}
If $f(x)=O(g(x))$, 
then there are positive constants $c_1$ and $n_{0}^{\prime}$
such that 
$$0\leq f(n)\leq c_1\,g(n)\mbox{ for all }n \ge n_{0}^{\prime},$$
and if $g(x)=O(h(x))$, 
then there are positive constants $c_2$ and 
$n_{0}^{\prime\prime}$ such that 
$$0\leq g(n)\leq c_2\,h(n)\mbox{ for all }n \ge n_{0}^{\prime\prime}.$$
Set $n_0=\max(n_0^{\prime},n_0^{\prime\prime})$ and $c_3=c_1\,c_2$.
Then
$$0\leq f(n)\leq c_1\,g(n)\leq c_1\,c_2\,h(n)=c_3\,h(n) 
\mbox{ for all }n \ge n_{0}.$$
Thus $f(x)=O(h(x))$.
\end{answer}
\end{exer}

%%%%%%%%%%%%%%%%%%%%%%%%%%%
\subsection{Proofs using limits}\label{subsec:limproof}
So far we have used the definitions of the various notations in all of our proofs. 
The following theorem provides another technique that is often much easier,
assuming you understand and are comfortable with limits.

\begin{thm}\label{thm:limits}
Let $f(n)$ and $g(n)$ be functions such that 
$$\lim_{n\rightarrow \infty}\frac{f(n)}{g(n)} = A.$$
Then
\begin{enumerate}
\item If $A=0$, then $f(n)=O(g(n))$, and $f(n)\neq \Theta(g(n))$.  That is, $f(n)=o(g(n))$.
\item If $A=\infty$, then $f(n)=\Omega(g(n))$, and $f(n)\neq \Theta(g(n))$. That is, $f(n)=\omega(g(n))$.
\item If $A\neq 0$ is finite, then $f(n)=\Theta(g(n))$. 
\end{enumerate}
\end{thm}

If the above limit does not exist, then you need to resort to
using the definitions or using some other technique.  
Luckily, in the analysis of algorithms the above approach works most of the time.

Before we see some examples, let's review a few limits you should know.

\begin{thm}\label{limitProps1}
Let $a$ and $c$ be real numbers.  Then
\begin{lenum}
\item $\displaystyle \lim_{n\rightarrow \infty} a = a$
\item If $a>0$, $\displaystyle \lim_{n\rightarrow \infty} n^a = \infty$
\item If $a<0$, $\displaystyle \lim_{n\rightarrow \infty} n^a = 0$
\item If $a>1$, $\displaystyle \lim_{n\rightarrow \infty} a^n = \infty$
\item If $0<a<1$, $\displaystyle \lim_{n\rightarrow \infty} a^n = 0$
\item If $c>0$, $\displaystyle \lim_{n\rightarrow \infty} \log_c n = \infty$.
\end{lenum}
\end{thm}
 
\begin{exa}
The following are examples based on Theorem~\ref{limitProps1}.
\begin{lenum}
\item
$\displaystyle\lim_{n\rightarrow\infty} 13 = 13$
\item
$\displaystyle\lim_{n\rightarrow\infty} n = \infty$
\item
$\displaystyle\lim_{n\rightarrow\infty} n^4 = \infty$
\item
$\displaystyle\lim_{n\rightarrow\infty} n^{1/2} = \infty$
\item
$\displaystyle\lim_{n\rightarrow\infty} n^{-2} = 0$
\item
$\displaystyle\lim_{n\rightarrow\infty} \left(\frac{1}{2}\right)^{n} = 0$
\item
$\displaystyle\lim_{n\rightarrow\infty} 2^{n} = \infty$
\item
$\displaystyle\lim_{n\rightarrow\infty} \log_2 n = \infty$
\end{lenum}
\end{exa}
 
Now it's your turn to try a few.
 
\begin{exer}
Evaluate the following limits
\begin{lenum}
  \setlength{\itemsep}{.25in}
\item
$\displaystyle\lim_{n\rightarrow\infty} \log_{10} n =$
\item
$\displaystyle\lim_{n\rightarrow\infty} n^3 =$
\item
$\displaystyle\lim_{n\rightarrow\infty} 3^n =$
\item
$\displaystyle\lim_{n\rightarrow\infty} \left(\frac{3}{2}\right)^n =$
\item
$\displaystyle\lim_{n\rightarrow\infty} \left(\frac{2}{3}\right)^n =$
\item
$\displaystyle\lim_{n\rightarrow\infty} n^{-1} =$
\item
$\displaystyle\lim_{n\rightarrow\infty} 8675309 =$
\end{lenum}
\begin{answer}
(a) $\infty$ 
(b) $\infty$ 
(c) $\infty$ 
(d) $\infty$ 
(e) $0$
(f) $0$
(g) $8675309$
\end{answer}
\end{exer}
 
\begin{exa}
Prove that  $5n^8=\Theta(n^8)$ using Theorem~\ref{thm:limits}.
\begin{sol}
Notice that 
$$\lim_{n\rightarrow\infty}\frac{5n^8}{n^8}= \lim_{n\rightarrow\infty}5=5,$$
so $f(n)=\Theta(n^8)$ by Theorem~\ref{thm:limits} (case 3).
\end{sol}
\end{exa}
 
The following theorem often comes in handy when using Theorem~\ref{thm:limits}.
\begin{thm}\label{thm:oneoverlimit}
If $\displaystyle \lim_{n\rightarrow \infty} f(n)=\infty$, then  
$\displaystyle \lim_{n\rightarrow \infty} \frac{1}{f(n)}=0$.
\end{thm}


\begin{exa}
Prove that  $n^2=o(n^4)$ using Theorem~\ref{thm:limits}.
\begin{sol}
Notice that 
$$\lim_{n\rightarrow\infty}\frac{n^2}{n^4}=
\lim_{n\rightarrow\infty}\frac{1}{n^2}=0,$$
so $f(n)=o(n^4)$ by Theorem~\ref{thm:limits} (case 1).
\end{sol}
\end{exa}

\begin{ques}
The proof in the previous example used Theorems~\ref{limitProps1}
and~\ref{thm:oneoverlimit}. How and where?

\noindent
\answerlinetwo
\begin{answer}
Theorem~\ref{limitProps1} part (b) implies that  
$\displaystyle\lim_{n\rightarrow\infty} n^2=\infty$. 
Since the limit being computed was actually 
$\displaystyle \lim_{n\rightarrow\infty}\frac{1}{n^2}$,
Theorem~\ref{thm:oneoverlimit}
was used to obtain the final answer of $0$ for the limit.
\end{answer}
\end{ques}

\begin{exer}
Prove that $3x^3=\Omega(x^2)$ using Theorem~\ref{thm:limits}.  Which case did you use?

\vspace*{1.75in}
\begin{answer}
Notice that $\lim_{x\rightarrow\infty}\frac{3x^3}{x^2}=
\lim_{x\rightarrow\infty}3x=\infty,$
so $3x^3=\omega(x^2)$ by the second case of the Theorem~\ref{thm:limits}, which also implies
that $3x^3=\Omega(x^2)$.
\end{answer}
\end{exer}

Here are a few more useful properties of limits.  
Read carefully.  These do not apply in all situations.

\begin{thm}\label{limitProps2}
Let $a$ be a finite real number and let $\displaystyle\lim_{n\rightarrow \infty} f(n)=A$ and 
$\displaystyle\lim_{n\rightarrow \infty} g(n)=B$, where $A$ and $B$ are finite real numbers. 
Then
\begin{lenum}
\item 
$\displaystyle \lim_{n\rightarrow \infty} a\,f(n) 
%=a\lim_{n\rightarrow \infty} f(n)
=a\,A$
\item 
$\displaystyle \lim_{n\rightarrow \infty} f(n)\pm g(n) 
%=\lim_{n\rightarrow \infty} f(n)\pm\lim_{n\rightarrow \infty} f(n)
=A\pm B$
\item 
$\displaystyle \lim_{n\rightarrow \infty} f(n)g(n) 
%=\lim_{n\rightarrow \infty} f(n)\lim_{n\rightarrow \infty} g(n)
=AB$
\item 
If $B\neq 0$, $\displaystyle \lim_{n\rightarrow \infty} \frac{f(n)}{g(n)} 
%=\frac{\lim_{n\rightarrow \infty}f(n)}{\lim_{n\rightarrow \infty} g(n)}
=\frac{A}{B}$
\end{lenum}
\end{thm}

We usually use the results from the previous theorem without explicitly mentioning them.

\begin{exa}
Find a tight bound on $f(x)=x^8+7x^7-10x^5-2x^4+3x^2-17$ using Theorem~\ref{thm:limits}.

\medskip

\noindent
{\bf Solution:}
We guess (or know, if we remember the solution to Example~\ref{exa:boundExample}) 
that $f(x)=\Theta(x^8)$.
To prove this, notice that 
\begin{eqnarray*}
\lim_{x\rightarrow\infty}
{x^8{+}7x^7{-}10x^5{-}2x^4{+}3x^2{-}\frac{17}{ x^8}}
& =& \lim_{x\rightarrow\infty}
\frac{x^8}{ x^8}+ \frac{7x^7}{ x^8}- \frac{10x^5}{ x^8}- \frac{2x^4}{ x^8}+
\frac{3x^2}{ x^8}- \frac{17}{ x^8}\\
& =& \lim_{x\rightarrow\infty}
 1+ \frac{7}{ x}- \frac{10}{ x^3}- \frac{2}{ x^4}+\frac {3}{ x^6}- \frac{17}{ x^8}\\
&=&1+ 0 - 0 - 0 + 0 - 0=1
\end{eqnarray*}
Thus, $f(x)=\Theta(x^8)$ by the Theorem~\ref{thm:limits}.
\end{exa}

Compare the proof above with the proof given in Example~\ref{exa:boundExample}.
It should be pretty obvious that using Theorem~\ref{thm:limits} makes the proof
a lot easier.  
Let's see another example that lets us compare the two proof methods.

\newpage % Another awkward break.

\begin{exa}
Prove that $f(x)=x^4-23 x^3 + 12x^2 + 15 x - 21=\Theta(x^4)$.
\begin{quote} 
{\bf Proof \#1}

\noindent
We will use the definition of $\Theta$.
It is clear that when $x\geq 1$, 
$$x^4-23 x^3 + 12x^2 + 15 x - 21\leq x^4+12 x^2+15 x
\leq x^4+12 x^4+15 x^4= 28 x^4.$$
Also, if $x\geq 88$, then $\frac{1}{2}x^4\geq 44x^3$ or 
$-44x^3\geq -\frac{1}{2}x^4$,
so we have that
$$x^4-23 x^3 + 12x^2 + 15 x - 21\geq x^4-23 x^3 - 21\geq 
x^4 - 23 x^3 - 21 x^3= x^4 - 44 x^3\geq \frac{1}{2}x^4.$$
%whenever 
%$$\frac{1}{2}x^4\geq 44x^3 \Leftrightarrow x\geq 88.$$
Thus
$$\frac{1}{2}x^4\leq x^4-23 x^3 + 12x^2 + 15 x - 21\leq 28 x^4,
\mbox{ for all } x\geq 88.$$
We have shown that $f(x)=x^4-23 x^3 + 12x^2 + 15 x - 21=\Theta(x^4).$
\hfill $\square$

\noindent
If you did not follow the steps in this first proof, you should really
review your algebra rules.

\noindent
{\bf Proof \#2}

\noindent
Since
\begin{eqnarray*}
\lim_{x\rightarrow\infty}\frac{
x^4-23 x^3 + 12x^2 + 15 x - 21}{x^4}
& =& \lim_{x\rightarrow\infty}
\frac{x^4}{ x^4} -\frac{23x^3}{ x^4} + \frac{12x^2}{ x^4} 
+ \frac{15x}{ x^4} - \frac{21}{ x^4}\\
& =& \lim_{x\rightarrow\infty}
1 - \frac{23}{ x}+ \frac{12}{ x^2} + \frac{15}{ x^3} - \frac{21}{ x^4}\\
&=& \lim_{x\rightarrow\infty}
1 - 0 + 0 + 0 - 0=1,
\end{eqnarray*}
$f(x)=x^4-23 x^3 + 12x^2 + 15 x - 21=\Theta(x^4)$
\hfill $\square$
\end{quote}
\end{exa}

\begin{exa}
Prove that $n(n+1)/2=O(n^3)$ using Theorem~\ref{thm:limits}.

\begin{pf}
Because 
${\displaystyle\lim_{n\rightarrow\infty}\frac{n(n+1)/2}{n^3} = 
\lim_{n\rightarrow\infty}\frac{n^2+n}{2n^3} = 
\lim_{n\rightarrow\infty}\frac{1}{2n} +\frac{1}{2n^2}= 0+0=0}$,
$n(n+1)/2=o(n^3)$, which implies that $n(n+1)/2=O(n^3)$.
\end{pf}
\end{exa}

\begin{exer}
Prove that $n(n+1)/2=\Theta(n^2)$ using Theorem~\ref{thm:limits}.

\vspace*{1.7in}
\begin{answer}
Notice that 
${\displaystyle \lim_{n\rightarrow\infty}\frac{n(n+1)/2}{n^2} = 
\lim_{n\rightarrow\infty}\frac{n^2+n}{2n^2} = 
\lim_{n\rightarrow\infty}\frac{1}{2} +\frac{1}{2n}= \frac{1}{2}+0=\frac{1}{2}}$,
so $n(n+1)/2=\Theta(n^2)$.
\end{answer}
\end{exer}


\begin{exer}
Prove that $2^x=O(3^x)$
\begin{lenum}
\item Using Theorem~\ref{thm:limits}.

\vspace*{1.5in}
\item Using the definition of Big-O.

\vspace*{1.5in}
\end{lenum}
\begin{answer}
(a)
Since 
${\displaystyle \lim_{x\rightarrow\infty} \frac{2^x}{ 3^x}
= \lim_{x\rightarrow\infty} \left(\frac{2}{ 3}\right)^x
= 0}$, the result follows.

(b)
If $x\geq 1$, then clearly $(3/2)^x\geq 1$, so 
$2^x\leq 
2^x\left(\frac{3}{ 2}\right)^x=
\left(\frac{2\times 3}{ 2}\right)^x=3^x.$
Therefore, $2^x=O(3^x)$.
\end{answer}
\end{exer}



Now is probably a good time to recall a very useful theorem for computing
limits, called {\bf l'Hopital's Rule}.  The version presented here is
restricted to limits where the variable
approaches infinity since those are the only limits of interest in our context.

\begin{thm}[l'Hopital's Rule]\label{thm:lhop}\index{l'Hopital's Rule}
Let $f(x)$ and $g(x)$ be differentiable functions.
If 
$$\lim_{x\rightarrow\infty}f(x)=\lim_{x\rightarrow\infty}g(x)=0 \mbox{ or}$$ 
$$\lim_{x\rightarrow\infty}f(x)=\lim_{x\rightarrow\infty}g(x)=\infty,$$ 
then
$$
\lim_{x\rightarrow\infty} \frac{f(x)}{ g(x)}=
\lim_{x\rightarrow\infty} \frac{f^{\prime}(x)}{ g^{\prime}(x)}$$
\end{thm}

%\newpage % A weird one. It's giving a bad break warning, but puts this on the next page even without the newpage.
  % What's the problem?

\begin{exa}
Since $\displaystyle \lim_{x\rightarrow\infty} 3x=\infty$ and
$\displaystyle \lim_{x\rightarrow\infty} x^2=\infty$, 
\begin{eqnarray*}
\lim_{x\rightarrow\infty} \frac{3x}{x^2}
&=&\lim_{x\rightarrow\infty} \frac{3}{2x} \, \, \, \,   \mbox{(l'Hopital)}\\
&=&\frac{3}{2}\lim_{x\rightarrow\infty} \frac{1}{x}\\
&=&\frac{3}{2}\,0\\
&=&0.
\end{eqnarray*}
\end{exa}

\begin{exa}
Since $\displaystyle \lim_{x\rightarrow\infty} 3x^2+4x-9=\infty$ and
$\displaystyle \lim_{x\rightarrow\infty} 12x=\infty$, 
\begin{eqnarray*}
\lim_{x\rightarrow\infty} \frac{3x^2+4x-9}{12x}
&=&\lim_{x\rightarrow\infty} \frac{6x+4}{12} \, \, \, \,   \mbox{(l'Hopital)}\\
&=&\lim_{x\rightarrow\infty} \frac{1}{2}x+\frac{1}{3}\\
&=&\infty
\end{eqnarray*}
\end{exa}

Now let's apply it to proving asymptotic bounds.

\begin{exa}
Show that $\log x=O(x)$.
\begin{pf}
Notice that 
\begin{eqnarray*}
\lim_{x\rightarrow\infty} \frac{\log x}{ x}
&=&\lim_{x\rightarrow\infty} \frac{\frac{1}{ x}}{ 1}\, \, \, \,  \mbox{(l'Hopital)} \\
&=&\lim_{x\rightarrow\infty} \frac{1}{ x}= 0.
\end{eqnarray*}
Therefore, $\log x=O(x)$.
\end{pf}
We should mention that applying l'Hopital's Rule in the first step is legal since 
$$\lim_{x\rightarrow\infty} \log x = \lim_{x\rightarrow\infty}  x=\infty.$$
\end{exa}

\begin{exa}
Prove that $x^3=O(2^x)$.
\begin{pf}
Notice that \begin{eqnarray*}
\lim_{x\rightarrow\infty} \frac{x^3}{2^x}
&=&\lim_{x\rightarrow\infty} \frac{3x^2}{2^x\ln(2)}\, \, \, \,  \mbox{(l'Hopital)} \\
&=&\lim_{x\rightarrow\infty} \frac{6x}{2^x\ln^2(2)}\, \, \, \,  \mbox{(l'Hopital)} \\
&=&\lim_{x\rightarrow\infty} \frac{6}{2^x\ln^3(2)}\, \, \, \,  \mbox{(l'Hopital)} \\
&=&0.
\end{eqnarray*}
Therefore, $x^3=O(2^x)$.

As in the previous example, at each step we checked that the functions on 
both the top and bottom
go to infinity as $n$ goes to infinity before applying l'Hopital's Rule.
Notice that we did not apply it in the final step since $6$ does not go to infinity.
\end{pf}
\end{exa}

\newpage

\begin{eval}
Prove that $7^x$ is an upper bound for $5^x$, but that it is not a tight bound.
\begin{spfenum}
\item
This is true if and only if $7^x$ always grows faster than $5^x$ 
which means $7^x-5^x>0$ for all $x\neq 0$.
If it is a tight bound, then $7^x-5^x=0$, which is only true for $x=0$.
So $7^x$ is an upper bound on $5^x$, but not a tight bound.

\noindent
\evallinetwo
\item 
${\displaystyle \lim_{x\rightarrow \infty}\frac{5^x}{7^x} = 
\lim_{x\rightarrow \infty}\frac{x\log 5}{x\log 7}.}$
Both go to infinity, but $x\log 7$ gets there faster, showing that $5^x=O(7^x)$.

\noindent
\evallinetwo
\item
${\displaystyle \lim_{x\rightarrow \infty}\frac{7^x}{5^x} = 
\lim_{x\rightarrow \infty}\left(\frac{7}{5}\right)^x=\infty}$
since $7/5>1$.  Thus $5^x=O(7^x)$ by the limit theorem.

\noindent
\evallinetwo
\end{spfenum}
\begin{answer}
\begin{pfevalenum}
\item
$7^x$ grows faster than $5^x$ does not mean $7^x-5^x>0$ for all $x\neq 0$.
For one thing, 
we are really only concerned about positive values of $x$.  
Further, we are specifically concerned about very large values of $x$.  
In other words, we want something to be true for all $x$ that are `large enough'.
Also, this statement does not take into account constant factors.  
Similarly, a tight bound does {\em not} imply that $7^x-5^x=0$.  
The bottom line:  This one is way off.  They are not conveying an understanding of
what `upper bound' really means, and they certainly haven't proven anything.
Frankly, I don't think they have a clue what they are trying to say in this proof.
\item
This one has several problems.  First, the application of l'Hopital's rule is incorrect.
The result should be $\displaystyle\lim_{x\rightarrow \infty}\frac{5^x\log 5}{7^x\log 7}$, 
which should make it obvious that l'Hopital's rule doesn't actually help in this case.
(The key to this one is to do a little algebra.)
The next problem is the statement `but $x\log 7$ gets there faster'.  
What exactly does that mean?
Asymptotically faster, or just faster?  If the former, it needs to be proven.  If the latter, that isn't enough to prove relative growth rates.  Finally, even if this showed that $5^x=O(7^x)$, that only shows
that $7^x$ is an upper bound on $5^x$.  It does not show that the bound is not tight.
The bottom line is that bad algebra combined with vague statements falls way short of
a correct proof.
\item
This proof is {\em very} close to being correct.  The main problem is that they only stated
that $5^x=O(7^x)$, but they also needed to show that $5^x\neq\Theta(7^x)$.  It turns out that the
theorem they mention also gives them that. So all they needed to add is `and $5^x\neq \Theta(7^x)$' at the end.
Technically, there is another problem---they should have taken the limit of $5^x/7^x$.  What they really showed
using the limit theorem is that $7^x=\omega(5^x)$, which is equivalent to $5^x=o(7^x)$.  It isn't a major problem,
but technically the limit theorem does not directly give them the result they say it does.  
If you are trying to prove that $f(x)$ is bounded by $g(x)$, put $f(x)$ on the top and $g(x)$ on the bottom.
\end{pfevalenum}
\end{answer}
\end{eval}
We should mention that it is important to remember to verify that l'Hopital's Rule applies
before just blindly taking derivatives. You can actually get the incorrect answer if
you apply it when it should not be applied.

\begin{exa}
Find and prove a simple tight bound for  $\sqrt{5n^2-4n+12}$.
\begin{sol}
We will show that $\sqrt{5n^2-4n+12}=\Theta(n)$.
Since we are letting $n$ go to infinity, we can assume that $n>0$.  In this case, 
$n=\sqrt{n^2}$.
Using this, we can see that 
$$\lim_{n\rightarrow\infty}\frac{\sqrt{5n^2-4n+12}}{n}
=\lim_{n\rightarrow\infty}\sqrt{\frac{5n^2-4n+12}{n^2}}
=\lim_{n\rightarrow\infty}\sqrt{5 - \frac{4}{n}+\frac{12}{n^2}}
=\sqrt{5}.
$$
Therefore, $\sqrt{5n^2-4n+12}=\Theta(n)$.
\end{sol}
\end{exa}

\newpage

\begin{exer}
Find and prove a good simple upper bound on $n\ln (n^2+1) + n^2\ln n$.
\begin{lenum}
\item Using the definition of Big-O.

\vspace*{3.5in}
\item Using Theorem~\ref{thm:limits}. 
You will probably need to use l'Hopital's Rule a few times.

\vspace*{4in}
\end{lenum}
\begin{answer}
You should have come up with $n^2\log n$ for the upper bound.  If you didn't,
now that you know the answer, go back and try to write the proofs before
reading them here.
(a) 
If $n> 1$, 
$$\ln (n^2+1) \leq \ln (n^2+n^2)= \ln (2 n^2) =
( \ln 2+ \ln n^2) \leq ( \ln n+ 2\ln n) =3\ln n$$
Thus when $n > 1$,
$$ n\ln (n^2+1) + n^2\ln n\leq n 3\ln n+n^2\ln n\leq
3 n^2 \ln n+n^2\ln n\leq  4 n^2 \ln n. $$
Thus, $n\ln (n^2+1) + n^2\ln n=O(n^2 \ln n).$
(You may have different algebra in your proof.  Just make certain
that however you did it that it is correct.)

(b)
$
\begin{array}[t]{llll}
\displaystyle\lim_{x\rightarrow\infty} \dfrac{{n\ln (n^2+1) + n^2\ln n}}{ n^2\ln n}
& =& \displaystyle\lim_{x\rightarrow\infty}
\dfrac{n\ln (n^2+1)}{ n^2\ln n}  + 1\\[.125in]
& =& \displaystyle1+ \lim_{x\rightarrow\infty}
\dfrac{\ln (n^2+1)}{ n\ln n} \\[.125in]
& =& \displaystyle1+ \lim_{x\rightarrow\infty}
\dfrac{\dfrac{2n}{n^2+1}}{1\cdot\ln n+n\cdot\frac{1}{n}} &\mbox{(l'Hopital)}\\[.125in]
& =& \displaystyle1+ \lim_{x\rightarrow\infty}
\dfrac{2n}{(n^2+1)(\ln n+1)} \\[.125in]
& =&\displaystyle 1+ \lim_{x\rightarrow\infty}
\dfrac{2}{2n(\ln n + 1) + (n^2+1)\cdot\frac{1}{n}} &\mbox{(l'Hopital)}\\[.125in]
& =&\displaystyle 1+ \lim_{x\rightarrow\infty}
\dfrac{2}{2n(\ln n + 1) + n+\frac{1}{n}}\\
&=&1+0=1.
\end{array}
$\\
Therefore,  $n\ln (n^2+1) + n^2\ln n=\Theta(n^2\log n)$.
\end{answer}
\end{exer}
\newpage

\begin{exa}
Find and prove a simple tight bound for  $n\log(n^2)+(n-1)^2\log(n/2)$.
\begin{sol}
First notice that 
$$n\log(n^2)+(n-1)^2\log(n/2)=2n\log n + (n-1)^2(\log n - \log 2).$$
We can see that this is $\Theta(n^2\log n)$ since 
\begin{eqnarray*}
\lim_{n\rightarrow\infty}\frac{n\log(n^2)+(n-1)^2\log(n/2)}{n^2\log n}
&=&\lim_{n\rightarrow\infty}\frac{2n\log n + (n-1)^2(\log n - \log 2)}{n^2\log n}\\
&=&\lim_{n\rightarrow\infty}\frac{2}{n} + \frac{(n-1)^2}{n^2}\frac{(\log n - \log 2)}{\log n}\\
&=&\lim_{n\rightarrow\infty}\frac{2}{n} + \left(1-\frac{1}{n}\right)^2\left(1-\frac{\log 2}{\log n}\right)\\
&=&0+(1-0)^2(1-0) = 1.
\end{eqnarray*}
\end{sol}
\end{exa}

\begin{exer}
Find and prove a simple tight bound for $(n^2-1)^5$.
You may use either the formal definition of $\Theta$ or Theorem~\ref{thm:limits}.
(The solution uses Theorem~\ref{thm:limits}.)

\vspace*{4in}
\begin{answer}
We can see that $(n^2-1)^5=\Theta(n^{10})$ since 
$$\lim_{n\rightarrow\infty}\frac{(n^2-1)^5}{n^{10}}
=\lim_{n\rightarrow\infty}\left(\frac{n^2-1}{n^2}\right)^5
=\lim_{n\rightarrow\infty}\left(1 - \frac{1}{n^2}\right)^5
=1
.$$
\end{answer}
\end{exer}

\newpage

\begin{exer}
Find and prove a simple tight bound for $2^{n+1}+5^{n-1}$.
You may use either the formal definition of $\Theta$ or Theorem~\ref{thm:limits}.
(The solution uses Theorem~\ref{thm:limits}.)

\vspace*{4in}
\begin{answer}
The following limit shows that $2^{n+1}+5^{n-1}=\Theta(5^n)$.
$$\lim_{n\rightarrow\infty}\frac{2^{n+1}+5^{n-1}}{5^n}
=\lim_{n\rightarrow\infty}\frac{2^{n+1}}{5^n}+\frac{5^{n-1}}{5^n}
=\lim_{n\rightarrow\infty}2\left(\frac{2}{5}\right)^n+\frac{1}{5}
=0+\frac{1}{5}.
$$
Note that we could also have shown that $2^{n+1}+5^{n-1}=\Theta(5^{n-1})$,
but that is not as simple of a function.
\end{answer}
\end{exer}

\newpage

\section{Common Growth Rates}\label{sec:growrates}
In this section we will take a look at the relative growth rates of various functions.

\noindent
\begin{figure}[h]
\begin{minipage}{.5\textwidth}
Figure~\ref{fig:comp1} shows the value of several functions for various values
of $n$ to give you an idea of their relative rates of growth.  
The bottom of the table is labeled relative to the last column 
so you can get a sense of how slow $\log_2 m$ and $\log_2(\log_2 m)$ grow.
For instance, the final row is showing that 
$\log_2(262144) = 18$ and 
$\log_2(\log_2 (262144) ) = 4.170$.

Figures~\ref{fig:comp2} and~\ref{fig:comp3} demonstrate that as $n$ increases, 
the constants and lower-order terms do not matter.
For instance, notice that although $100n$ is much larger than $2^n$ for small values of $n$,
as $n$ increases, $2^n$ quickly gets much larger than $100n$. 
Similarly, in Figure \ref{fig:comp3}, notice that when $n=74$, $n^3$ and $n^3+234$ are virtually
the same.

\vspace*{.35in}
\end{minipage}
\hfill
\begin{minipage}{.47\textwidth}
{\footnotesize
$
%\begin{array}{|r|r|r|r|r|r|}
%\hline
%\log n&n&n\log n&n^2&n^3&2^n\rule{0in}{.15in}\\
%\hline
%0&1&0&1&1&2\\
%0.6931&2&1.39&4&8&4\\
%1.099&3&3.30&9&27&8\\
%1.386&4&5.55&16&64&16\\
%1.609&5&8.05&25&125&32\\
%1.792&6&10.75&36&216&64\\
%1.946&7&13.62&49&343&128\\
%2.079&8&16.64&64&512&256\\
%2.197&9&19.78&81&729&512\\
%2.303&10&23.03&100&1000&1024\\
%2.398&11&26.38&121&1331&2048\\
%2.485&12&29.82&144&1728&4096\\
%2.565&13&33.34&169&2197&8192\\
%2.639&14&36.95&196&2744&16384\\
%2.708&15&40.62&225&3375&32768\\
%2.773&16&44.36&256&4096&65536\\
%2.833&17&48.16&289&4913&131072\\
%2.890&18&52.03&324&5832&262144\\
%\hline
%\log\log m&\log m&&&&m\\
%\hline
%\end{array}
\begin{array}{|r|r|r|r|r|r|}
\hline
\log_2 n&n&n\ln n&n^2&n^3&2^n\rule{0in}{.15in}\\
\hline
0.000&1&0&1&1&2\\
1.000&2&1.39&4&8&4\\
1.585&3&3.30&9&27&8\\
2.000&4&5.55&16&64&16\\
2.321&5&8.05&25&125&32\\
2.585&6&10.75&36&216&64\\
2.807&7&13.62&49&343&128\\
3.000&8&16.64&64&512&256\\
3.170&9&19.78&81&729&512\\
3.321&10&23.03&100&1000&1024\\
3.460&11&26.38&121&1331&2048\\
3.585&12&29.82&144&1728&4096\\
3.700&13&33.34&169&2197&8192\\
3.807&14&36.95&196&2744&16384\\
3.907&15&40.62&225&3375&32768\\
4.000&16&44.36&256&4096&65536\\
4.087&17&48.16&289&4913&131072\\
4.170&18&52.03&324&5832&262144\\
\hline
\log_2\log_2 m&\log_2 m&&&&m\\
\hline
\end{array}
$
}\caption{A comparison of growth rates}\label{fig:comp1}
%\captionof{figure}{A comparison of growth rates}\label{fig:comp1}
\end{minipage}
\end{figure}




\begin{figure}[h]
\begin{minipage}{.4\textwidth}
{\footnotesize
$
\begin{array}{|r|r|r|r|r|r|}
\hline
n&100 n&n^2&11 n^2&n^3&2^n\rule{0in}{.15in}\\
\hline
1&100&1&11&1&2\\
2&200&4&44&8&4\\
3&300&9&99&27&8\\
4&400&16&176&64&16\\
5&500&25&275&125&32\\
6&600&36&396&216&64\\
7&700&49&539&343&128\\
8&800&64&704&512&256\\
9&900&81&891&729&512\\
10&1000&100&1100&1000&1024\\
11&1100&121&1331&1331&2048\\
12&1200&144&1584&1728&4096\\
13&1300&169&1859&2197&8192\\
14&1400&196&2156&2744&16384\\
15&1500&225&2475&3375&32768\\
16&1600&256&2816&4096&65536\\
17&1700&289&3179&4913&131072\\
18&1800&324&3564&5832&262144\\
19&1900&361&3971&6859&524288\\
\hline
\end{array}
$
}
\caption{Constants don't matter}\label{fig:comp2}
\end{minipage}
\hfill
\begin{minipage}{.5\textwidth}
{\footnotesize
$
\begin{array}{|r|r|r|r|r|r|}
\hline
n&n^2&n^2-n&n^2+99&n^3&n^3+234\rule{0in}{.15in}\\ % the rule gives some needed vertical space
\hline
2&4&2&103&8&242\\
6&36&30&135&216&450\\
10&100&90&199&1000&1234\\
14&196&182&295&2744&2978\\
18&324&306&423&5832&6066\\
22&484&462&583&10648&10882\\
26&676&650&775&17576&17810\\
30&900&870&999&27000&27234\\
34&1156&1122&1255&39304&39538\\
38&1444&1406&1543&54872&55106\\
42&1764&1722&1863&74088&74322\\
46&2116&2070&2215&97336&97570\\
50&2500&2450&2599&125000&125234\\
54&2916&2862&3015&157464&157698\\
58&3364&3306&3463&195112&195346\\
62&3844&3782&3943&238328&238562\\
66&4356&4290&4455&287496&287730\\
70&4900&4830&4999&343000&343234\\
74&5476&5402&5575&405224&405458\\
\hline
\end{array}
$
}
\caption{Lower-order terms don't matter}\label{fig:comp3}
\end{minipage}
\end{figure}

Figures \ref{fig:growthgraph1} through \ref{fig:growthgraph5} give a graphical
representation of relative growth rates of functions. In these diagrams,
\verb+**+ means exponentiation.  For instance, 
\verb+x**2+ means $x^2$.

%% Got rid of most undefull hboxes by replacing 
%\begin{minipage}{.47\textwidth} for each figure with just multicols.
\begin{figure}[ht]
\begin{center}
\begin{multicols}{2}
\sputfig{fig/plot2.eps}{.4}
\caption{Slow growing functions.}\label{fig:growthgraph1}
\columnbreak
\sputfig{fig/plot1.eps}{.4}
\caption{Polynomials.}
\end{multicols}
\end{center}

\smallskip
\begin{center}
\begin{multicols}{2}
\sputfig{fig/plot3.eps}{.4}
\caption{Polynomials and an exponential.  
It looks like $x^4$ grows faster than $2^x$, but see Fig~\ref{fig:f2}.}\label{fig:f1}
\columnbreak
\sputfig{fig/plot4.eps}{.4}
\caption{Polynomials and an exponential with larger $n$. Clearly $2^n$ grows faster than $n^4$.}\label{fig:f2}
\end{multicols}
\end{center}

\smallskip
\begin{multicols}{2}
\sputfig{fig/plot5.eps}{.4}
\caption{\raggedright Notice that as $n$ gets larger, the constants eventually matter less.}
\label{fig:growthgraph5}
\vspace*{0in} % gets rid of \vbox underfull warning.
\columnbreak
\vspace*{.5in}
\end{multicols}
\end{figure}

It is important to point out that you should {\em never} rely on the graphs of functions
to determine relative growth rates.  That is the point of Figures~\ref{fig:f1} and~\ref{fig:f2}.
Although graphs sometimes give you an accurate picture of the relative growth rates of
the functions, they might just as well present a distorted view of the data 
depending on the values that are used on the axes.  Instead, you should use the
techniques we develop in this section.

Next we present some of the most important results about the relative growth rate of some 
common functions.  We will ask you to prove each of them. 
Theorems \ref{thm:limits} and \ref{thm:lhop} will help you do so.
You will notice that most of the theorems are using little-o, not Big-O.
Hopefully you understand the difference.  If not, review those definitions before
continuing.

We begin with something that is pretty intuitive: higher powers grow faster than lower powers.

\begin{thm}\label{thm:polyComp}
Let $a<b$ be real numbers.  Then 
$n^a=o(n^b)$.
\end{thm}

\begin{exa}
According to Theorem~\ref{thm:polyComp}, $n^2=o(n^3)$ and $n^5=o(n^{5.1})$.
\end{exa}

\begin{exer}\label{exer:polyComp}
Prove Theorem~\ref{thm:polyComp}.
(Hint: Use Theorem~\ref{thm:limits} 
and do a little algebra before you try to compute the limit.)

\vspace*{2in}
\begin{answer}
Since $a<b$, $b-a>0$.  
Therefore, $\displaystyle \lim_{n\rightarrow\infty} \frac{n^a}{n^b} 
=  \lim_{n\rightarrow\infty} n^{a-b} =  \lim_{n\rightarrow\infty} \frac{1}{n^{b-a}}=0$.
By Theorem~\ref{thm:limits}, $n^a=o(n^b)$.
\end{answer}
\end{exer}


The next theorem tells us that exponentials with different bases do not grow at the same rate.
More specifically, the higher the base, the faster the growth rate.

\begin{thm}\label{thm:expComp}
Let $0<a<b$ be real numbers.  Then 
$a^n=o(b^n)$.
\end{thm}

\begin{exa}
According to Theorem~\ref{thm:expComp}, $2^n=o(5^n)$ and $4^n=o(4.5^n)$.
\end{exa}

\begin{exer}
Prove Theorem~\ref{thm:expComp}.
(See the hint for Exercise~\ref{exer:polyComp}.)

\vspace*{2in}
\begin{answer}
Since $a<b$, $a/b<1$.  
Therefore, $\displaystyle \lim_{n\rightarrow\infty} \frac{a^n}{b^n} 
=  \lim_{n\rightarrow\infty}\left(\frac{a}{b}\right)^n = 0$.
By Theorem~\ref{thm:limits}, $a^n=o(b^n)$.
\end{answer}
\end{exer}

%Stupid page break issues.  I can't get the previous box to size to the bottom of the page.
%It either goes to the next page or the next paragraph goes on this page.  Not sure why.
%I should eventually figure out what this box environment is doing.  CAC, 6/10/15.

\newpage

Recall that a logarithmic function is the inverse of an exponential function.
That is, $b^{x} = n$ is equivalent to $x = \log_{b} n$.
The following identity is very useful.

%\begin{thm}\label{thm:rellogs}
%Let $a$ and $b$ be positive real numbers.  Then 
%$$\log_a b = \frac{\log_a c}{\log_b c}.$$
%This can be rearranged as 
%$$\frac{\log_a c}{\log_a b} =\log_b c 
%\mbox{\, \, \,  or \, \, \, \, }
%(\log_a b)(\log_b c)=log_a c.$$
%\end{thm}
\begin{thm}\label{thm:rellogs}
Let $a$, $b$, and $x$ be positive real numbers with $a\neq 1$ and $b\neq 1$.  Then 
$$\log_a x = \frac{\log_b x}{\log_b a}.$$
\end{thm}

\begin{exa}
Most calculators can compute $\ln n$ or $\log_{10} n$, but are unable to compute
logarithms with any given base.  But Theorem~\ref{thm:rellogs} allows you to do so.
For instance, you can compute $\log_2 39$ as $\log_{10} 39/\log_{10} 2$.
\end{exa}

Notice that the formula in Theorem~\ref{thm:rellogs}  can be rearranged as 
$(\log_b a)(\log_a x) = \log_b x.$
This form should make it evident that 
changing the base of a logarithm just changes the value
by a constant amount.  This leads to the following result.
\begin{cor}\label{cor:thetalogs}
Let $a$ and $b$ be positive real numbers with $a\neq 1$ and $b\neq 1$.

\noindent
Then $\log_a n=\Theta(\log_b n)$.
\begin{pf}
Follows from the definition of $\Theta$ and
Theorem~\ref{thm:rellogs}.
\end{pf}
\end{cor}

\begin{exa}
According to Corollary~\ref{cor:thetalogs}, $\log_2 n =\Theta(\log_{10}n)$ 
and $\ln n =\Theta(\log_2 n)$.
\end{exa}

Corollary~\ref{cor:thetalogs} is stating that all logarithms have the same 
rate of growth regardless of their bases.
That is, the base of a logarithm does not matter when it is used
in asymptotic notation.
Because of this, the base is often omitted  in asymptotic notation.  
In computer science, it is usually safe to assume that 
the base of logarithms is $2$ if it is not specified.

\begin{exer}
Indicate whether each of the following is {\em true} (T) or {\em false} (F).

\vspace*{-2ex}
\begin{lenum}
\itemtf
$2^n=\Theta(3^n)$
\itemtf
$2^n=o(3^n)$
\itemtf
$3^n=O(2^n)$
\itemtf
$\log_3 n=\Theta(\log_2 n)$
\itemtf
$\log_2 n=O(\log_3 n)$
\itemtf
$\log_{10} n=o(\log_3 n)$
\end{lenum}
\begin{answer}
(a) False since $3^n$ grows faster than $2^n$.
(b) True since $2^n$ grows slower than $3^n$.
(c) False since $3^n$ grows faster than $2^n$, which means it does not grow
slower or at the same rate.
(d) True since they both have the same growth rate.  Remember, exponentials with
different bases have different growth rates, but logarithms with different bases
have the same growth rate.
(e) True since they have the same growth rate.  Remember that if $f(n)=\Theta(g(n))$, 
then $f(n)=O(g(n))$ and $f(n)=\Omega(g(n))$.
(f) False since they have the same growth rate, so $\log_{10} n$ does not grow slower
than $\log_3 n$.
\end{answer}
\end{exer}
Next we see that logarithms grow slower than positive powers of $n$.  


\begin{thm}\label{thm:logpower1}
Let $b>0$ and $c>1$ be real numbers. Then $\log_c(n)=o(n^b)$.
\end{thm}

\begin{exa}
According to Theorem~\ref{thm:logpower1}, $\log_2 n= o(n^2)$,  $\log_{10} n= o(n^{1.01})$, 
and $\ln n =o(\sqrt{n})$.
\end{exa}

\begin{exer}
Prove Theorem~\ref{thm:logpower1}.
(Hint: This is easy if you use Theorems~\ref{thm:limits} and~\ref{thm:lhop})

\vspace*{2.5in}
\begin{answer}
Using l'Hopital's rule, we have 
$\displaystyle \lim_{n\rightarrow\infty} \frac{\log_c(n)}{n^b} 
=  \lim_{n\rightarrow\infty}\frac{\frac{1}{n\ln(c)}}{b\, n^{b-1}}
=  \lim_{n\rightarrow\infty}\frac{1}{\ln(c)b\, n^b}
=0$ since $b>0$.  Thus, Theorem~\ref{thm:limits} tells us that $\log_c n=o(n^b)$.
\end{answer}
\end{exer}

More generally, the next theorem states that 
any positive power of a logarithm grows slower than any positive
power of $n$.  Since this one is a little tricky, we will provide the proof.
In case you have not seen this notation before, you should know that $\log^a n$
means $(\log n)^a$, which is {\em not} the same thing as $\log (n^a)$.

\begin{thm}\label{thm:logpowernpower}
Let $a>0$, $b>0$, and $c>0$ be real numbers. Then $\log_c^a (n)=o(n^b)$.
In other words, any power of a log grows slower than any polynomial.
\begin{pf}
First, we need to know that if $a>0$ is a constant, and
$\displaystyle\lim_{n\rightarrow\infty}f(n)=C$, then
$$\displaystyle\lim_{n\rightarrow\infty}(f(n))^a
=\left(\displaystyle\lim_{n\rightarrow\infty}f(n)\right)^a
=C^a.$$
Using this and the limit computed in the proof of Theorem~\ref{thm:logpower1}, we have that 
$$\displaystyle \lim_{n\rightarrow\infty} \frac{\log_c^a(n)}{n^b} 
=\displaystyle \lim_{n\rightarrow\infty} \left(\frac{\log_c(n)}{n^{b/a}}\right)^a 
=\left(\displaystyle \lim_{n\rightarrow\infty}\frac{\log_c(n)}{n^{b/a}}\right)^a 
=0^a=0.
$$
Thus, Theorem~\ref{thm:limits} tells us that $\log_c^a(n)=o(n^b)$.
\end{pf}
\end{thm}

\begin{exa}
According to Theorem~\ref{thm:logpowernpower}, $\log_2^4 n= o(n^2)$,
$\ln^{10} n =o(\sqrt{n})$,
and

\noindent
$\log_{10}^{1,000,000} n= o(n^{.00000001})$.
\end{exa}

Finally, any exponential function with base larger than $1$ grows faster than any polynomial.


\begin{thm}\label{thm:polyexp}
Let $a>0$ and $b>1$ be real numbers. Then $n^a=o(b^n)$.
\end{thm}

\begin{exa}
According to Theorem~\ref{thm:polyexp}, it is easy to see that $n^2=o(2^n)$, $n^{15}=o(1.5^n)$,
and $n^{1,000,000}=o(1.0000001^n)$.
\end{exa}

There are several ways to prove Theorem~\ref{thm:polyexp}, including using repeated 
applications of l'Hopital's rule, using induction, or doing a little
algebraic manipulation and using one of several clever tricks.
But the techniques are beyond what we generally need in the course, so we will
omit a proof (and, perhaps more importantly, we will not ask you to provide a proof!).

\begin{fib}
Fill in the following blanks with  $\Theta$, $\Omega$, $O$, or $o$.
You should give the most precise answer possible. (e.g. If you put $O$, but
the correct answer is $o$, your answer is correct but not precise enough.)
\begin{lenum}
\item $n(n-1)=\blankline{.5}(500n^2)$.
\item $50n^2=\blankline{.5}(.001n^4)$.
\item $\log_2{n}=\blankline{.5}(\ln{n})$.
\item $\log_2\left(n^2\right)=\blankline{.5}(\log_2^2(n))$.
\item $2^{n-1}=\blankline{.5}(2^n)$.
\item $5^{n}=\blankline{.5}(3^n)$.
\item $(n-1)!=\blankline{.5}(n!)$.
\item $n^3=\blankline{.5}(2^n)$.
\item $\log^{100}n=\blankline{.5}(1.01^n)$.
\item $\log^{100}n=\blankline{.5}(n^{1.01})$.
\end{lenum}
\begin{answer}
(a) $\Theta$;
(b) $o$ ($O$ is correct, but not precise enough.);
(c) $\Theta$;
(d) $o$ ($O$ is correct, but not precise enough.);
(e) $\Theta$ since $2^n=2\,2^{n-1}$;
(f) $\Omega$ (Technically it is $\omega$, but I'll let it slide if you put $\Omega$ since
 we haven't used $\omega$ much.);
(g) through (j) are all $o$ ($O$ is correct, but not precise enough.)
\end{answer}
\end{fib}

An alternative notation for little-o is $\ll$.  In other words, $f(n)=o(g(n))$ iff
$f(n)\ll g(n)$.  This notation is useful in certain contexts, including the following
comparison of the growth rate of common functions.  The previous theorems in this section
provide proofs of some of these relationships.  The others are given without proof.

\begin{thm}
Here are some relationships between the growth rates of common functions:
$$c \ll \log n \ll log^{2} n \ll \sqrt n \ll 
n \ll n \log n \ll n^{1.1} \ll n^{2} \ll  n^3 \ll n^4 \ll 2^n\ll 3^n\ll n!\ll n^n$$
\end{thm}

You should convince yourself that each of the relationships given in the previous
theorem is correct.  You should also memorize them or (preferably) understand why
each one is correct so you can `recreate' the theorem.

\begin{exer}
Give a $\Theta$ bound for each of the following functions. You do not need to prove them.
\begin{lenum}
  \setlength{\itemsep}{.5in}
\item $f(n) = n^5 + n^3 + 1900 + n^7 + 21n + n^2 $
\item $f(n) = (n^2+23n+19 ) (n^2+23n+ n^3+19 ) n^3 $  (Don't make this one harder than it is)
\item $f(n) = n^2  + 10,000n + 100,000,000,000 $
\item $f(n) = 49*2^n+34*3^n $
\item $f(n) = 2^n + n^5 + n^3$
\item $f(n) = n \log n + n^2$
\item $f(n) = \log^{300} n + n^{.000001}$
\item $f(n) = n!\log n + n^n+3^n$
\end{lenum}
\vspace*{.45in}
\begin{answer}
If your answers do not all start with $\Theta$, go back and redo them before reading the answers.  Your answers should match the following {\em exactly}.
(a) $\Theta(n^7)$.
(b) $\Theta(n^8)$.
(c) $\Theta(n^2)$.
(d) $\Theta(3^n)$.
(e) $\Theta(2^n)$.
(f) $\Theta(n^2)$.
(g) $\Theta(n^{.000001})$.
(h) $\Theta(n^n)$.
\end{answer}
\end{exer}

\newpage % awkward pagebreak

\begin{exer}
Rank the following functions in increasing rate of growth.
Clearly indicate if two or more functions have the same growth rate.
Assume the $\log$s are base 2.
\medskip

\noindent
$x$,\, \,  
$x^2$,\, \,  
$2^x$,\, \,  
$10000$,\, \,    
$\log^{300} x$,\, \,   
$x^{5}$,\, \,   
$\log x$,\, \,   
$x^{\log 3}$,\, \,    
$x^{.000001}$,\, \,     
$3^x$,\, \,  
$x \log(x)$,\, \,  
$\log (x^{300})$,\, \,   
$\log(2^x)$

\vspace*{3.5in}
\begin{answer}
Here is the correct ranking (where $\sim$ indicates two functions grow at the same rate):

\noindent
$10000$,\, 
$\log x$ $\sim$ $\log (x^{300})$,\, 
$\log^{300} x$,\, 
$x^{.000001}$, \,   
$x$ $\sim$ $\log(2^x)$,  \,   
$x \log(x)$,  \,   
$x^{log_2 3}$, \,   
$x^2$, \,   
$x^{5}$, \,   
$2^x$,  \,   
$3^x$.
\end{answer}
\end{exer}



\newpage
%%%%%%%%%%%%%%%%%%%%%%%%%%%%%%%%%%%%%%%%%%%%%%%%%%%%%%%%
\section{Algorithm Analysis}\label{sec:algAnalysis}

The overall goal of this chapter is to deal with a seemingly simple question:
{\em Given an algorithm, how good is it?}
I say ``seemingly'' simple because unless we define what we mean by ``good'', we cannot
answer the question. 
Do we mean how elegant it is?  How easy it is to understand?
How easy it is to update if/when necessary?  Whether or not it can be generalized?

Although all of these may be important questions, in algorithm analysis we are usually more interested
in the following two questions:  
{\em How long does the algorithm take to run, and how much space (memory) 
does the algorithm require.}  
In fact, we follow the tradition of most books and focus our discussion on the first question.
This is usually reasonable since the amount of memory used by most algorithms is not large 
enough to matter.  There are times, however, when analyzing the space required by an algorithm is
important.  For instance, when the data is really large (e.g. the graph that represents 
friendships on Facebook) or when you are implementing a {\em space-time-tradeoff} algorithm.

Although we have simplified the question, we still need to be more specific.  
What do we mean by ``time''?
Do we mean how long it takes in real time (often called {\em wall-clock time})?
\index{wall-clock time}  
Or the actual amount of time our processor used (called {\em CPU time})?
\index{CPU time}  
Or the exact {\em number of instructions} (or {\em number of operations}) executed?  

\begin{ques}
Why aren't wall-clock time and CPU time the same? 

\noindent
\answerlinethree
\begin{answer}
Modern computers use multitasking to perform several tasks (seemingly) at the same time.
Therefore, if an algorithm takes 1 minute of real time (wall-clock time), it might be that
58 seconds of that time was spent running the algorithm, but it could also be the case that
only 30 seconds of that time were spent on that algorithm, and the other 30 seconds spent
on other processes.  In this case, the CPU time would be 30 seconds, but the wall-clock time
60 seconds.

Further complicating matters is increasing availability of machines
with multiple processors.  If an algorithm runs on 4 processors rather than one, it 
might take 1/4th the time in terms of wall-clock time, but it will probably take the
same amount of CPU time (or close to it).
\end{answer}
\end{ques}

Because the running time of an algorithm is greatly affected by the 
characteristics of the computer system 
(e.g. processor speed, number of processors, amount of memory, file-system type, etc.),
the running time does not necessarily provide a comparable measure, 
regardless of whether you use CPU time or wall-clock time.
The next question asks you to think about why.

\begin{ques}
Sue and Stu were competing to write the fastest algorithm to solve a problem.
After a week, Sue informs Stu that her program took 1 hour to run.  
Stu declared himself victorious since his program took only 3 minutes. 
But the real question is this:  Whose algorithm was more efficient?  
Can we be certain Stu's algorithm was better than Sue's?  Explain.
(Hint:  Make sure you don't jump to any conclusion too quickly.  Think about all of 
the possibilities.)

\noindent
\answerlinethree
\begin{answer}
We cannot be certain whose algorithm is better with the given information.  
Maybe Sue used a {\em TRS-80 Model IV} from the 1980s to run her program and Stu used 
{\em Tianhe-2} (The fastest computer in the world from about 2013-2014).  
In this case, it is possible that if Sue ran her program on {\em Tianhe-2} 
it would have only taken 2 minutes, making her the real winner.
\end{answer}
\end{ques}

The answer to the previous question should make it clear that you cannot compare the running times of algorithms
if they were run on different machines.  
Even if two algorithms are run on the same computer, the wall-clock times may not be comparable.

\begin{ques}
Why isn't the {\em wall-clock time} of two algorithms that are run on the same computer always
a reliable indicator of their relative performances?

\noindent
\answerlinethree
\begin{answer}
As has already been mentioned, other processes on the machine can have a significant influence
on the wall-clock time.  For instance, if I run two CPU-intensive programs at once, the wall-clock time of each might be about twice what it would be if I ran them one at a time.
If they are run on a machine with multiple cores
the wall-clock time might be closer to the CPU-time.  But other processes that are 
running can still throw off the numbers.
\end{answer}
\end{ques}

In fact, if you run the same algorithm on the same machine multiple times, 
it will not always take the same amount of time.  
Sometimes the differences between trial runs can be significant.

\begin{ques}
If two algorithms are run on the same machine, can we reliably compare the {\em CPU-times}?

\noindent
\answerlinethree
\begin{answer}
For the most part, yes.  This is especially true if the running times of the algorithms
are not too close to each other (in other words, if one of the algorithms is significantly
faster than the other).
However, the number of other processes running on the machine 
can have an influence on CPU-time.  For instance, if there are more processes running,
there are more context switches, and depending on how the CPU-time is counted, these context
switches can influence the runtime.  So although comparing the CPU-time of two algorithms
that are run on the same computer gives a pretty good indication of which is better, it
is still not perfect.
\end{answer}
\end{ques}

So the CPU-time turns out to be a pretty good measure of algorithm performance. 
Unfortunately, it does not really allow one to compare two {\em algorithms}.  It only allows
us to compare specific {\em implementations of the algorithms}.  
It also requires us to implement the algorithm in an actual programming language before we
even know how good the algorithm is (that is, before we know if we should even spend the
time to implement it).

But we can analyze and compare algorithms before they are implemented if we use
the {\em number of instructions} as our measure of performance.  
There is still a problem with this measure.  
What is meant by an ``instruction''?  
When you write a program in a language such as Java or C++, it is not executed exactly
as you wrote it.  It is compiled into some sort of machine language.  
The process of compiling does not
generally involve a one-to-one mapping of instructions, 
so counting Java instructions versus C++ instructions
wouldn't necessarily be fair.  
On the other hand, we certainly do not want to look at the machine code in
order to count instructions---machine code is ugly.
Further, when analyzing an algorithm, should we even take into account the exact 
implementation in a particular language, or should we analyze the algorithm apart from implementation?

O.K., that's enough of the complications.  Let's get to the bottom line.  When analyzing algorithms, we 
generally want to ignore what sort of machine it will run on and what language it will be implemented in.
We also generally do not want to know {\em exactly} how many instructions it will take.  
Instead, we want to know the {\em rate of growth} of the number of instructions.  
This is sometimes called the {\em asymptotic running time} of an algorithm.
In other words, as the size of the input increases, how does that affect the number of instructions executed?
We will typically use the notation from Section~\ref{sec:asymptotic} 
to specify the running time of an algorithm.
We will call this the {\em time complexity} (or often just {\em complexity})
of the algorithm.

\subsection{Analyzing Algorithms}

Given an algorithm, the {\em size of the input} is exactly what it sounds like---the amount of space
required to specify the input.  For instance, if an algorithm operates on an array of size $n$, 
we generally say the input is of size $n$.
For a graph, it is usually the number of vertices or the number of vertices and edges. 
When the input is a single number, things get more complicated for reasons I
do not want to get into right now.  We usually don't need to worry about this, though.

Algorithm analysis involves determining the size of the input, $n$,  and 
then finding a function based on $n$ that tells us how long the algorithm will take if the
input is of size $n$.  By ``how long'' we of course mean how many operations.

\begin{exa}[Sequential Search]\index{sequential search}
\index{search!sequential}
Given an array of $n$ elements, often one needs to determine if a given number $val$ is in the array.  
One way to do this is with the {\em sequential search} algorithm that simply looks through all of 
the elements in the array until it finds it or reaches the end.  The most common version of this
algorithm returns the index of the element, or $-1$ if the element is not in the array.
Here is one implementation.
\begin{code}
    int sequentialSearch(int a[],int n, int val) {
        for(int i=0;i<a.size();i++) {
            if(a[i]==val) {
                return i;
            }
        }
        return -1;
    }
\end{code}
What is the size of the input to this algorithm?
\begin{sol}
There are a few possible answers to this question.  The input technically consists of an array
of $n$ elements, the numbers $n$, and the value we are searching for.  So we could consider
the size of the input to be $n+2$.  However, typically we ignore constants with input sizes.
So we will say the size of the input is $n$.
\end{sol}
\end{exa}

In general, if an algorithm takes as input an array of size $n$ and some constant number of
other numeric parameters, we will consider the size of the input to be $n$.  

\begin{exer}
Consider an algorithm that takes as input an
$n$ by $m$ matrix, an integer $v$, and a real number $r$.  What is the size of the input?

\noindent
\answerline
\begin{answer}
This one is a little more tricky.  
The answer is $n\cdot m$ since this is how many entries are in the matrix.  
Sometimes we need to use two numbers to specify the input size.
As suggested previously, we will ignore the size of the two other pieces of data.
\end{answer}
\end{exer}

\begin{exa}
How many operations does \scode{sequentialSearch} take on an array of size $n$?
\begin{sol}
As mentioned above, we consider $n$ as the size of the input.  
Assigning $i=0$ takes one instruction.  
Each iteration through the \verb+for+ loop increments $i$,
compares $i$ with $a.size()$, and compares $a[i]$ with $val$.  
Don't forget that accessing $a[i]$ and calling \scode{a.size()} 
each take (at least) one instruction. 
Finally, it takes an
instruction to return the value.
If the $val$ is in the array at position $k$, the algorithm will take $2+5k=\Theta(k)$
operations, the $2$ coming from the assignment \scode{i=0} and the return statement.  
If  $val$ is not in the array, the algorithm 
takes $2+5n=\Theta(n)$ instructions.
\end{sol}
\end{exa}

This last example should bring up a few questions.  Did we miss any instructions?  
Did we miss any possible outcomes that would give us a different answer?  How exactly should we specify our analysis?

Let's deal with the possible outcomes question first.  Generally speaking, when we analyze an algorithm we want to know what happens in one of three cases:  
The  best case, the average case, or the worst case. 
When thinking about these cases, 
we always consider them for a given value of $n$ (the input size).  
We will see in a moment why this matters.
 
As the name suggests, when performing a {\em best case} analysis, 
we are trying to determine the smallest
possible number of instructions an algorithm will take.  
Typically, this is the least useful type of analysis.
If you have experienced a situation when someone said something like 
``it will only take an hour (or a day) to fix your cell phone,'' 
and it actually took 3 hours (or days), you will understand why.  

When determining the best-case performance of an algorithm, remember that we need
to determine the best-case performance for a given input size $n$.  This is important
since otherwise every algorithm would take a constant amount of time in the best case simply by giving it an input of the smallest possible size (typically $0$ or $1$).  That sort of
analysis is not very informative.

\begin{rem}
When you are asked to do a best-case analysis of an algorithm, remember that it is
implied that what is being asked is the best-case analysis {\bf for an input of size $n$}.
This actually applies to average and worst-case analysis as well, but it is easier to
make this mistake when doing a best-case analysis.
\end{rem}

{\em Worst case} analysis considers what is the largest number of instructions that will execute
(again, for a given input size $n$). 
This is probably the most common analysis, and typically the most useful.  
When you pay Amazon for guaranteed
2-day delivery, you are paying for them to guarantee a worst-case delivery time.  
However, this analogy is imperfect.  
When you do a worst-case analysis, you know the algorithm will {\em never} take longer
than what your analysis specified, but occasionally an Amazon delivery is lost or delayed.
When you perform a worst-case analysis of an algorithm, you always consider what can happen
that will make an algorithm take as long as possible, so it will never take longer than the worst-case analysis implies.
 
The {\em average case} is a little more complicated, both to define and to compute.  
The first problem is determining what ``average'' means for a particular input and/or algorithm.  
For instance, what does an ``average'' array of values look like?
The second problem is that even with a good definition, computing the average case
complexity is usually much more difficult than the other two.
It also must be used appropriately.  
If you know what the average number of instructions for an algorithm is, 
you need to remember that sometimes it might take less time and 
sometimes it might take more time--possibly significantly more time.

I sometimes use the term {\em expected running time} instead of one of these three.
This is almost synonymous with {\em average case}, 
but when I use this term I am being less formal.  
Thus, I will not necessarily do a complete
average case analysis to determine what I call the expected running time.  
Think of it as being how long an algorithm will {\em usually} take.  It will typically
coincide with either the average- or worst-case complexity.

\begin{exa}
Continuing the  \scode{sequentialSearch} example, notice that our analysis above reveals that the
best-case performance is $7=\Theta(1)$ operations 
(if the element sought is the first one in the array) and the worst-case performance is
$2+5n=\Theta(n)$ operations (if the element is not in the array).  
If we assume that the element we are searching for is
equally likely to be anywhere in the array or not in the array, 
then the average-case performance 
should be about $2+5(n/2)=\Theta(n)$ operations.  
We will do a more thorough average-case analysis of this algorithm shortly.
\end{exa}

Notice that in the previous example, 
the average- and worst-case complexities are the same.  This makes sense. 
We estimate that the average case takes about half as long as the worst case.  But no matter how
large $n$ gets, it is still just half as long.  
That is, the {\em rate of growth} of the average and worst-case running times are the same.
Also note the logic we used to obtain the best-case complexity of $\Theta(1)$.
We did not say the best case was $\Theta(1)$ because the best-case input was 
an array of size one.  Instead it is $\Theta(1)$ because in the best case the element we
are searching for is the first element of the array, no matter how large the array is.

Here is another important question: 
How do we know we counted all of the operations?
As it turns out, we don't actually care.
This is good because determining the exact number is very difficult, if not impossible. 
Recall that we said we wanted to know the rate of growth of an algorithm, not the exact number of instructions. 
As long as we count all of the ``important'' ones, we will get the correct rate of growth.  
But what are the ``important'' ones?
The term {\em abstract operation} is sometimes used to describe the operations that we will count.
Typically you choose one type of operation or a set of operations that you know will be performed 
the most often and consider those as the abstract operation(s).

\begin{exa}\index{sequential search}\index{search!sequential}
The analysis of \scode{sequentialSearch} can be done more easily than in the
previous example.  
We repeat the algorithm here for convenience.
\begin{code}
    int sequentialSearch(int a[],int n, int val) {
        for(int i=0;i<a.size();i++) {
            if(a[i]==val) {
                return i;
            }
        }
        return -1;
    }
\end{code}
Notice that the comparison (\scode{a[i]==val}) is executed as often as any other instruction.
Therefore if we count the number of times that instruction executes, we can use that 
to determine the rate of growth of the running time.
  
In the best case the comparison is executed once (if the element being searched for is
the first one in the array), so the best-case complexity is $\Theta(1)$.

In the worst case the comparison is executed $n=\Theta(n)$ times
(if the element being searched for is either at the end or not present in the array).
  
As before, we expect the average case to be about $n/2=\Theta(n)$, although
in the next example we will do a more complete analysis.

Notice that we obtained the same answers here as we did above when we tried to 
take into account every operation.
\end{exa}

\begin{exa}
Let's determine the average-case complexity of \scode{sequentialSearch}.
If we assume that the element is equally likely to be anywhere in the array,
then there is a $1/n$ chance that it will be in any given spot.  If it is in the
first spot, the comparison executes once.  If it is in the second spot, it executes twice.
In general it takes $k$ comparisons if it is in the $k$th spot.
Since each possibility has a $1/n$ chance, the average expected search time is
$$\sum_{k=1}^{n}\frac{k}{n}=\frac{1}{n}\sum_{k=1}^{n}k=\frac{1}{n}\frac{n(n+1)}{2}
=\frac{n+1}{2}=\Theta(n).$$

Our analysis simplified things a bit---we didn't take into account the possibility that
the element was not in the array.  To do so, let's assume the element searched for is
equally likely to be anywhere in the array or not in the array.
That is, there is now a $1/(n+1)$ chance that it will be in any of the $n$ spots in the
array and a $1/(n+1)$ chance that it is not in the array. (We divide by $n+1$
because there are now $n+1$ possibilities, each equally likely.)  If it is not in the array,
the number of comparisons is $n$. 
In this case the expected time would be 
$$\sum_{k=1}^{n}\left(\frac{k}{n+1}\right) + \frac{n}{n+1}
=\frac{1}{n+1}\left(\sum_{k=1}^{n}k + n\right)
=\frac{1}{n+1}\left(\frac{n(n+1)}{2}+ n\right)
=\frac{n^2+3n}{2(n+1)}
=\Theta(n).$$
We'll leave it to you to prove that 
$\frac{n^2+3n}{2(n+1)} =\Theta(n)$. (Use Theorem~\ref{thm:limits} and a little algebra).
\end{exa}

The previous example demonstrates how performing an average-case analysis is typically
much more difficult than the other two, even with a relatively simple algorithm.  
In fact, did we even do it correctly?
Is it a valid assumption that there is a $1/(n+1)$ chance that the element searched for
is not in the array?  If we are searching for a lot of values in a small array, 
perhaps it is the case that most of the values we are searching for are {\em not} in the array.  
Maybe it is more realistic to assume there is a $50\%$ chance it is in the array 
and $50\%$ chance that it is not in the array.  
I could propose several other reasonable
assumptions, too.  As stated before, it can be difficult to define ``average.''
In this case it actually doesn't matter a whole lot because under any reasonable
assumptions the average-case analysis will always come out as $\Theta(n)$.

As you might be able to imagine, things get much more complicated as the algorithms
get more complex.  This is one of the reasons that in some cases we will skip or
gloss over the details of the average-case analysis of an algorithm.

It is important to make sure that you choose the operation(s) you will count carefully so
your analysis is correct.  
In addition, you need to look at every instruction in the algorithm to determine whether
or not it can be accomplished in constant time.  If some step takes
longer than constant time, that needs to be properly taken into consideration.
In particular, consider function/method calls and operations on data structures 
very carefully.  For instance, if you see a method call like \verb+insert(x)+
or \verb+get(x)+, you cannot just assume they take constant time.  You need to determine
how much time they actually take.

\begin{rem}
When you are asked for the {\em complexity} of an algorithm, you should 
do the following three things:
\begin{enumerate}
\item
Give the best, average, and worst-case complexities unless otherwise specified.
Sometimes the average case is quite complicated and can be skipped.
\item
Give answers  in the form of $\Theta(f(n))$ for some function $f(n)$,
or $O(f(n))$ if a tight bound is not possible.  The function $f(n)$ you choose should
be as simple as possible.  For instance, instead of $\Theta(3n^2+2n+89)$, you should use
$\Theta(n^2)$ since the constants and lower order terms don't matter.
\item
Clearly justify your answers by explaining how you arrived at them in sufficient detail.
\end{enumerate}
\end{rem}

\newpage

\begin{exa}
What is the complexity of \verb+max(x,y)+?  Justify your answers.
\begin{code}
	int max(int x, int y) {
	    if(x >= y) {
	        return x;
	    } else {
	        return y;
	    }
	}
\end{code}
\begin{sol}
No matter what, the algorithm does a single comparison followed by a return statement.  Therefore, in the best, average, and worst case, \verb+max+ takes
about $2$ operations. Therefore, the complexity is always $\Theta(1)$ 
(otherwise known as {\em constant}).
\end{sol}
\end{exa}


\begin{exer}
Analyze the following algorithm that finds the maximum value in an array.
Start by deciding which operation(s) should be counted.
Don't forget to give the best, worst, and average-case complexities.
\begin{code}
    int maximum(int a[],int n) {
        int max = int.MIN_VAL; 
        for (int i=0; i<n; i++)
            max = max(max, a[i]);
        return max; 
    } 
\end{code}
\vspace*{3in}

\begin{answer} 
We focus on the assignment (=) inside the loop and ignore the other instructions.
This should be fine since assignment occurs at least as often as any other instruction.  
In addition, it is important to note that \verb+max+ takes constant time
(did you remember to explicitly say this?), 
as do all of the other operations, so we aren't under-counting.
It isn't too difficult to see that the assignment will occur $n$ times for an array of size $n$
since the code goes through a loop with $i=0,\ldots, n-1$.
Thus, the complexity of \scode{maximum} is always $\Theta(n)$.  That is, $\Theta(n)$ is the best, average, and worst-case complexity of \scode{maximum}.
\end{answer}
\end{exer}

When an algorithm has no conditional statements (like the \verb+maximum+ algorithm from the previous exercise),
or at least none that can cause the algorithm to end earlier, the best, average, and worst-case
complexities will usually be the same.  I say usually because there is always the possibility
of a weird algorithm that I haven't thought of that could be an exception.

\newpage

\begin{exa}\label{exa:sequentialloops}
Give the complexity of the following code.
\begin{code}
     int q=0;
     for (int i=1; i<=n; i++) { 
         q=q+i*i;
     }
     for (int j=1; j<=n; j++) {
         q=q*j;
     }
\end{code}
\begin{sol}
This algorithm has two independent loops, 
each of which do slightly different things. 
Thus, we cannot pick a single operation to count.  
Instead we will pick the assignment statements that involve $q$.
That is, we will use both \verb#q=q+i*i# and \verb+q=q*j+.
The first assignment executes $n$ times since the first loop executes for
every value of $i$ from $1$ to $n$.  The second loop also executes its assignment
$n$ times for the same reason.  Since the loops happen one after another, we
add the number of operations, so the total is $n+n=2n$ assignment statements.
Since there are no conditional statements, this is the best, worst, and
average-case number of assignment statements.  Thus, the complexity for all three
cases is $\Theta(n)$.
 
\end{sol}
\end{exa}
 

\begin{exa}\label{exa:nestedloops}
Give the complexity of the following code.
\begin{code}
    double V = 0;
     for (int i=1; i<=n; i++) { 
         for (int j=1; j<=n; j++) { 
             V=V+A[i]*A[j];
         }
     }
\end{code}
\begin{sol}
Clearly the assignment (\scode{V=A[i]*A[j])} occurs the most often.
The inner loop\footnote{{\em Always} analyze from the inside out.  The more practice you get, the
more it will be obvious that this is the only way that will consistently work.}
always executes $n$ times, each time doing one assignment.
The outer loop executes $n$ times, and each time it executes, it executes the inner loop.
Therefore the total time is $n\cdot n=\Theta(n^2)$.  This is the best, worst, and average
case complexity since nothing about the input can change what the algorithm does.

Here is another way to think about it. 
The inner loop executes the assignment statement $n$ times every time it executes.
The first time through the outer loop, 
the whole inner loop executes an calls the assignment $n$ times.
The second time through the outer loop, 
the whole inner loop executes an calls the assignment $n$ times.
This happens all the way until the $n$th time through the outer loop during which the
whole inner loop executes an calls the assignment $n$ times.
Thus, the total number of times the assignment is called is 
$n+n+\cdots+n$ times (where there are $n$
terms in the sum), which is just $n\cdot n$.  Thus the complexity is $\Theta(n^2)$.
 
\end{sol}
\end{exa}
 
Sometimes people mistakenly think the algorithm Example~\ref{exa:sequentialloops} 
takes $\Theta(n^2)$ operations.  id
But it is not executing one loop inside another loop.  
It is executing one loop $n$ times followed by another loop $n$ times.
On the other hand, the algorithm in Example~\ref{exa:nestedloops}
does {\em not} take $n+n$ operations.  
It is not executing one loop $n$ times followed by another loop $n$ times.
It is executing one loop $n$ times, and each of those $n$ times it is executing another
loop that takes $n$ time.

Here is an analogy.  If you climb a flight of $10$ stairs followed by another flight of $10$ 
stairs, you climbed a total of $10+10=20$ stairs.
Now assume you go into a building that has $10$ floors.  There are $10$ steps between
floors (so it takes $10$ steps to get from floor 1 to 2, etc.)  If you climb to the 
top of the building, how many stairs did you climb? It is $10+10+\cdots+10$ (where there are
$10$ terms in the sum), which is $100=10^2$.  How does this relate to the previous examples?
Simple.  In the first case, you executed:
\begin{code}
    for(stair 1 through 10)
        climb stair
    for(stair 1 through 10)
        climb stair
\end{code}
and in the second case you executed:
\begin{code}
    for(floors 1 through 10)
        for(stair 1 through 10)
            climb stair
\end{code}
Do you see the resemblance to the code from 
Examples~\ref{exa:sequentialloops} and~\ref{exa:nestedloops}?
And do you see how we are really performing the same analysis?
 
It is important to be careful not to jump to conclusions when analyzing algorithms.  
For instance, a double-nested for-loop should always take $\Theta(n^2)$ to execute, right?
\begin{exer}
What is the worst-case complexity of the following algorithm?
\begin{code}
    int k=50;
    for (i = 0; i < n; i ++) {
        for (j = 0; j < k; j ++) {
            a[i][j] = b[i][j] * x;
        }
    }
\end{code} 
\vspace*{2in}
\begin{answer}
The line in the inner for loop takes constant time (let's call it $c$).
The inner loop executes $k=50$ times, each time doing $c$ operations.
Thus the inner loop does $50\cdot c$ operations, which is still just a constant.
The outer loop executes $n$ times, each time executing the inner loop, which takes $50\cdot c$ operations.
Thus, the whole algorithm takes $50\cdot c\cdot n=\Theta(n)$ time.
\end{answer}
\end{exer}

If you read the solution to the previous exercise 
(which you definitely should have---{\em always} read the solutions!), 
you will see that you need to be careful not to jump to conclusions too quickly.
A double-nested loop does not always mean an algorithm takes $\Theta(n^2)$ time.
But does it guarantee it will take $O(n^2)$ (in other words, no more than quadratic
time)?

\newpage

\begin{exer}
What is the worst-case complexity of the following algorithm?
\begin{code}
    for (i = 0; i < n; i ++) {
        for (j = 0; j < n*n; j ++) {
            a[i][j] = b[i][j] * x;
        }
    }
\end{code} 
\vspace*{2.25in}
\begin{answer}
The line in the inner for loop takes constant time (let's call it $c$).
The inner loop executes $n^2$ times since $j$ is going from
$0$ to $n^2-1$, so each time the inner loop executes, it does $cn^2$ operations.
The outer loop executes $n$ times, each time executing the inner loop.
Thus, the total time is $n\times cn^2=\Theta(n^3)$.

This is an example of an algorithm with a double-nested loop 
that is {\em worse than} $\Theta(n^2)$.  The point of this exercise is to make
it clear that you should never jump to conclusions too quickly when analyzing
algorithms.  Read the limits on loops very carefully!
\end{answer}
\end{exer}
%%%%%%%%%%%%%%%%%%%%%%%%%%%%%%%%%%%%%%%%%%%%%%%%%%%%%%%%
\subsection{Common Time Complexities}\label{section:growthRates}

We have already discussed the relative growth rates of functions.
In this section we apply that understanding to the analysis of algorithms.
That is, we will discuss common time complexities that are encountered
when analyzing algorithms. 
Let $n$ be the size of the input and $k$ a constant. 
We will briefly discuss each of the following complexity classes,
which are listed (mostly) in order of rate of growth.

\begin{itemize}
\item Constant: $\Theta(k)$, for example $\Theta(1)$
\item Logarithmic: $\Theta(\log_{k}n)$
\item Linear: $\Theta(n)$ 
\item $n\log n$: $\Theta(n \log_{k}n)$
\item Quadratic: $\Theta(n^{2})$ 
\item Polynomial: $\Theta(n^k)$
\item Exponential: $\Theta(k^{n})$
\end{itemize}

\begin{df}[Constant]\index{constant growth rate}\index{growth rate!constant}
An algorithm with running time $\Theta(1)$ 
(or $\Theta(k)$ for some constant $k$)
is said to have {\bf constant} complexity.  
Note that this does not necessarily mean that the algorithm takes 
exactly the same amount of time for all inputs, but it {\bf does}
mean that there is some number $K$ such that it always takes 
no more than $K$ operations.
\end{df}

\begin{exa}
The following algorithms have constant complexity.

\begin{minipage}[t]{.45\textwidth}\begin{code}
int FifthElement(int A[],int n) 
{
    return A[4];
}

\end{code} 
\end{minipage}
\hfill
\begin{minipage}[t]{.5\textwidth}\begin{code}
    int PartialSum(int A[],int n) {
        int sum=0;
        for(int i=0;i<42;i++) 
           sum=sum+A[i];
        return sum;
    }
\end{code} 
\end{minipage}
The algorithm \verb+FifthElement+ just indexes into an array and returns that value.
Since array indexing takes constant time, as does returning a single value,
this algorithm clearly takes just constant time, no matter how large $n$ is.

The algorithm \verb+PartialSum+ might seem to take $O(n)$ time since it contains
a loop.  But don't jump to conclusions too quickly.  Notice that the loop executes
$42$ times, regardless of how large $n$ might be.  All of the other operations
(both in and out of the loop) takes constant time.  Thus, the overall complexity
is something like $c_1+42*c_2$, where $c_1$ is the time it takes to do the operations
outside the loop, and $c_2$ is the time it takes to execute the code in the loop 
each time it executes, including the comparison and increment in the for loop itself.
Since both $c_1$ and $c_2$ are constant, so is $c_1+42*c_2$.  Thus, the algorithm
takes constant time.
\end{exa}

\begin{exer}
Which of the following algorithms have constant complexity?   
Briefly justify your answers.
\begin{lenum}
\item
The \verb+AreaTrapezoid+ algorithm from Example~\ref{exa:areaTrap}.

\noindent
\answerlinetwo
\item
The \verb+factorial+ algorithm from Example~\ref{exa:factorial}.

\noindent
\answerlinetwo
\item 
The \verb+absoluteValue+ algorithm from Problem~\ref{pro:ab1}.

\noindent
\answerlinetwo
\end{lenum}
\begin{answer}
(a) \verb+AreaTrapezoid+ is constant.  
(b) \verb+factorial+is not constant.  
It should be easy to see that it has a complexity of $\Theta(n)$.
(c)  \verb+absoluteValue+ is constant if we assume \verb+sqrt+ takes constant time.
\end{answer}
\end{exer}

\begin{df}[Logarithmic]\index{logarithmic growth rate}\index{growth rate!logarithmic}
Algorithms with running time $\Theta(\log n)$ are said to have {\bf logarithmic} complexity.
As the input size $n$ increases, so does the running time, but very slowly.
Logarithmic algorithms are typically found when the algorithm can systematically 
ignore fractions of the input.
\end{df}

\begin{exa}
In Example~\ref{exa:binSearchIter} we will see that 
{\em binary search} has complexity $\Theta(\log n)$. 
\end{exa}

\begin{df}[Linear]\index{linear growth rate}\index{growth rate!linear}
Algorithms with running time $\Theta(n)$ are said to have {\bf linear} complexity.
As $n$ increases, the run time increases in proportion with $n$.
Linear algorithms access each of their $n$ inputs at most some constant number of times.
\end{df}

\begin{exa}
The following are linear algorithms.
\medskip

\noindent
\begin{minipage}{.45\textwidth}
\begin{code}
  void sumFirstN(int n) {
      int sum=0; 
      for (int i=1;i<=n;i++)
      sum = sum + i;  
  }
\end{code}
\end{minipage}
\hfill
\begin{minipage}{.45\textwidth}
\begin{code}
  void mSumFirstN(int n) {
      int sum=0;
      for(int i=1;i<=n;i++)
          for(int k=1;k<7;k++)
              sum = sum + i;
  }
\end{code} 
\end{minipage}

It is pretty easy to see that \verb+sumFirstN+ takes linear time since it contains a
single for loop that executes $n$ times and does a constant amount of work each time.

At first glance it may seem that \verb+mSumFirstN+ takes $\Theta(n^2)$ time since it
has a double nested loop.  You will think about why it is actually $\Theta(n)$ in 
the next question.
\end{exa}

\begin{ques}
Why is the complexity of \verb+mSumFirstN+ from the previous example
$\Theta(n)$ and not $\Theta(n^2)$?

\noindent
\answerlinethree
\begin{answer}
Although it has a nested loop, the inside loop always executes $6$ times, which
is a constant.  So the algorithm takes about $6\cdot c\cdot n=\Theta(n)$ operations, not $\Theta(n^2)$.
\end{answer}
\end{ques}


\begin{df}[$n \log n$]
Many divide-and-conquer algorithms have complexity $\Theta(n\log n)$.
These algorithms break the input into a constant number
of subproblems of the same type, solve them independently, and then
combine the solutions together.  Not all divide-and-conquer algorithms
have this complexity, however.
\end{df}

\begin{exa}
Two of the most well known sorting algorithms, {\em Quicksort} and {\em merge sort}, 
have an average case complexity of $\Theta(n\log n)$.  
We will do a complete analysis of both algorithms in Chapter~\ref{ch:recurrences}.
\end{exa}

\begin{df}[Quadratic]\index{quadratic growth rate}\index{growth rate!quadratic}
Algorithms with a running time of $\Theta(n^2)$ are said to have {\bf quadratic} complexity.
As $n$ doubles, the running time quadruples.
\end{df}

\begin{exa}
The following algorithm is quadratic.
\begin{code}
    int compute_sums(int A[], int n) {
        int M[n][n];
        for (int i=0;i<n;i++) 
            for (int j=0;j<n;j++)
                M[i][j]=A[i]+A[j];
        return M;
    }
\end{code} 
This one is pretty easy to see since it has double nested loops that each
execute $n$ times, and the amount of work done in the inner loop is constant.
\end{exa}


\begin{exer}
Which of the following algorithms have quadratic complexity?   
Briefly justify your answers.
\begin{lenum}
\item
The \verb+factorial+ algorithm from Example~\ref{exa:factorial}.

\noindent
\answerlinetwo
\item
An algorithm that tries to find the smallest element in an array of size $n\times n$
by searching through the entire array.

\noindent
\answerlinetwo
\end{lenum}
\begin{answer}
(a) Since \verb+factorial+ has a complexity of $\Theta(n)$, it is not quadratic.
(b) Since there are $n^2$ entries it to consider, 
the algorithm takes $\Theta(n^2)$ time, so it would be quadratic.\footnote{Technically this is linear with respect to the {\em size} of the input since the
size of the input is $n^2$.  But it is quadratic in $n$.  In either case, it is
$\Theta(n^2)$.}
\end{answer}
\end{exer}


\begin{ques}
In a previous course you may have encountered several quadratic sorting algorithms.  Name them.
(Note: We will analyze two of them soon.)

\noindent
\answerline
\begin{answer}
Bubble sort, selection sort, and insertion sort are three of them that you may have seen before.
\end{answer}
\end{ques}

\begin{df}[Polynomial]\index{polynomial growth rate}
\index{growth rate!polynomial}
Algorithms with running time $\Theta(n^k)$ for some constant $k$ are 
said to have {\bf polynomial} complexity.  We call them 
{\bf polynomial-time algorithms}.  
Note that linear and quadratic are special cases of polynomial.
When we say an {\bf efficient}\index{efficient algorithm} 
algorithm exists to solve a problem,
we typically mean a polynomial-time algorithm.
\end{df}

\begin{exa}
As we will see in Example~\ref{exa:matmult}, \verb+MatrixMultiply+ takes $\Theta(n^3)$ time.
Since $3$ is a constant, that is a polynomial-time algorithm.
We will also mention Strassen's algorithm that has a complexity of about $\Theta(n^{2.8})$.
That is also a polynomial-time algorithm.  It's actual complexity is 
$\Theta(n^{\log_2 7})$.
\end{exa}


\begin{df}[Exponential]\index{exponential growth rate}\index{growth rate!exponential}
Algorithms with running time $\Theta(k^n)$ for some constant $k$ are 
said to have {\bf exponential} complexity. 
Since exponential algorithms can only be run for small values of $n$, 
they are not considered to be efficient.
Brute-force algorithms are often exponential.
\end{df}

\begin{exa}
Since there are $2^n$ binary numbers of length $n$, 
an algorithm that lists all binary numbers of length $n$ would take $\Theta(2^n)$ time,
which is exponential.
\end{exa}

\begin{rem}
As we have already seen, 
exponentials with different bases do not grow at the same rate.
Thus, two exponential algorithms do not belong to the same complexity class 
unless the base of the exponent is the same.
In other words, $a^n\neq\Theta(b^n)$ unless $a=b$.
\end{rem}

Let me end on a very important note regarding analysis of algorithms and asymptotic growth of functions.
If algorithm $A$ is faster than algorithm $B$, 
then the running time of $A$ is less than the running time of $B$.
On the other hand, if $A$'s running time is asymptotically faster than the running time of $B$,
that means $B$ is a faster algorithm!  
In other words, the words fast/slow need to be reversed when discussing
algorithm speeds versus the growth of the functions.
Put simply: {\em A faster growing complexity means a slower algorithm, and vice-versa}.

%%%%%%%%%%%%%%%%%%%%%%%%%%%%%%%%%%%%%%%%%
\subsection{Basic Sorting Algorithms}
In this section we will analyze algorithms that are slightly more complex than the previous examples. 
All of these algorithms take a list of $n$ integers and places them in ascending order. This is referred
to as {\em sorting} the lists. This idea of course can be expanded to any type of list as long as there
is a clear way to compare the elements to determine which is larger.
If you have previously taken a data structures course, you should be familiar with these algorithms.
In case you haven't had such a course, we provide very brief descriptions of the algorithms. 
But our focus here is on determining how efficient the algorithms are, not on exactly how the 
algorithms work.

\begin{exa}\label{exa:bubble}\index{bubble sort}\index{sort!bubble}
The \verb+bubblesort+ algorithm is one of the first sorting algorithms that students learn about.
Here is one implementation of the algorithm.
\begin{code}
    void bubblesort(int a[],int n) {
        for(int i=n-1;i>0;i--) {
            for(int j=0;j<i;j++) {
                if(a[j] > a[j+1]) {
                    swap(a,j,j+1);
                }
            }
        }
    }
    
\end{code}
If you have seen \verb+bubblesort+ previously, 
hopefully you know by now that you should {\em never} use it.
It is one of the worst sorting algorithms. 
Because of this, we won't even describe the algorithm here and just focus on analyzing it.\\

Find the complexity of \verb+bubblesort+, where $n$ is the size of the array $a$.
\begin{sol}
First, notice that the input size is $n$ since we are sorting an array 
with $n$ elements.

Example~\ref{exa:swapforArray} gives an implementation of \verb+swap+
that takes constant time (verify this!).  The conditional statement,
including the swap, takes constant time (we'll call it $c$, as usual), 
regardless of whether or not
the condition is true.  It takes longer if the condition is true, but
it is constant either way---about $3$ operations (array indexing ($\times 2$) and
comparison) versus about $6$ (the swap adds about $3$).

The inner loop goes from $j=0$ to $j=i-1$, 
so it executes $i$ times and takes $ci$ time.  
But what is $i$?  This is where things
get a little more complicated than in the previous examples.  
Notice that the outer loop is changing the value of $i$.  
We need to look at this a little more carefully.
\begin{enumerate}
\item 
The first time through the outer loop $i=n-1$.  
So the inner loop takes $c(n-1)$ time.
\item
The second time through the outer loop $i=n-2$.  
So the inner loop takes $c(n-2)$ time.
\item
The $k$th time through the outer loop $i=n-k$.  
So the inner loop takes $c(n-k)$ time.
\item
This goes all the way to the $n$th time through the outer loop when 
$i=1$ and the inner loop takes $c\cdot 1$ time.
\end{enumerate}
The outer loop is simply causing the inner loop to be executed over and over again,
but with different parameters (specifically, it is changing the limit on the inner loop).  
Thus, we need to add up the time taken for all of these calls to the inner loop.  
Doing so, we see that the total time required for \verb+bubblesort+ is 
\begin{eqnarray*}
c(n{-}1)+c(n{-}2)+c(n{-}3)+\cdots+c\, 1&=&c( (n{-}1) + (n{-}2) + (n{-}3)+\cdots + 1)\\
&=&c(1+2+3+\cdots+(n-1))\\
&=&c\sum_{k=1}^{n-1}k\\
&=&c\frac{(n-1)n}{2}\\
&=&\Theta(n^2)
\end{eqnarray*}
Thus, the complexity (worst, best, and average) of \verb+bubblesort+ is $\Theta(n^2)$.
\end{sol}
\end{exa}

\begin{rem}
Part way through our analysis of \verb+bubblesort+ we had $k$ as part of our complexity.
But notice that the $k$ did not show up as part of the final complexity.  This is because
in the context of the entire algorithm, $k$ has no meaning.  It is a local variable from the 
algorithm that we needed to use to determine the overall complexity of the algorithm.
{\bf The only variables that should appear in the complexity of an 
algorithm are those that are related to the size of the input.}
\end{rem}

\newpage
\begin{ques}
In the best case, the code in the conditional statement in \verb+bubblesort+ never 
executes.  Why does this still result in a complexity of $\Theta(n^2)$?

\noindent
\answerlinethree
\begin{answer}
As we mentioned in our analysis, executing the conditional statement takes about
$3$ operations, and if it is true, about $3$ additional operations are performed.  So
the worst case is no more than about twice as many operations as the best case.  In other
words, we are comparing $c\cdot n^2$ to $2c\cdot n^2$, both of which are $\Theta(n^2)$.
\end{answer}
\end{ques}

In reality, the best and worst case performance of \verb+bubblesort+ are different---the
worst case is about twice as many operations.  But when we are discussing the {\em complexity}
of algorithms, we care about the asymptotic behavior---that is, 
what happens as $n$ gets larger.  In that case, the difference is still just a factor of $2$.
The best and worst-case complexities have the same growth rate (quadratic).

Consider how this is different if the best-case complexity of an algorithm is $\Theta(n)$ 
and the worst-case complexity is $\Theta(n^2)$.  As $n$ gets larger, the gap between the
performance in the best and worst cases also gets larger.  In this case, 
the best and worst-case complexities are not the same since one is linear and the other
is quadratic.

\begin{rem}
If an algorithm contains nested loops and the limit on one or more of
the inner loops depends on a variable from an outer loop, analyzing the
algorithm will generally involve one or more summations, as it did with the previous example.
As mentioned previously, variables related to those loops that are 
used in your analysis (e.g. $i$, $j$, $k$, etc.) should never show up in your final answer!
They have no meaning in that context.
\end{rem}


\begin{exa}\index{insertion sort}\index{sort!insertion}
One of the best simple sorting algorithms is \verb+insertionSort+. 
It works much like one might sort a hand of playing cards.
The first card is in order since there is one one. Compare the second card to the first card and
swap if they are out of order. Now the first two are in order. Now pick up the third card and
go down the list and place it where it belongs. Just keep repeating this process until the whole
list is sorter. Here is an implementation of that idea:

\begin{code}
    void insertionSort(int a[], int n) {
        for (int i=1;i<n;i++) { 
            int v=a[i];  
            int j=i-1;
            while (j >= 0 && a[j] > v) {
                a[j+1] = a[j];  
                j--; 
            }
            a[j+1]=v;
        }
    }
\end{code}

Find the complexity of \verb+insertionSort+, where $n$ is the size of the array $a$.
\begin{sol}
The code inside the while loop takes constant time.
The loop can end for one of two reasons---if $j$ gets to $0$, or if \verb+a[j]>v+.  
In the worst case, it goes until $j=0$.
Since $j$ starts out being $i$ at the beginning, and it is decremented in the loop,
that means the loop executes $i$ times in the worst case.  

The for loop (the outer loop) 
changes the value of $i$ from $1$ to $n-1$, executing a constant amount
of code plus the while loop each time. So the $i$th time through the outer loop
takes $c_1+c_2i$ operations. 
We will simplify this to just $i$ operations---you 
can think of it as counting the number of assignments in the while loop if you wish.  
So the worst-case complexity is 
$$\sum_{i=1}^{n-1}i = \frac{(n-1)n}{2} =\Theta(n^2).$$
This happens, by the way, if the elements in the array start out in reverse order.

In the best case, the loop only executes once each time because \verb+a[j]>v+ is always true
(which happens if the array is already sorted).  
In this case, the complexity is $\Theta(n)$ since the outer loop executes $n-1$ times,
each time doing a constant amount of work.
\end{sol}

We should point out that if we had done our computations using $c_1+c_2i$ 
instead of $i$ we would have arrived at the same answer, 
but it would have been more work:
$$\sum_{i=1}^{n-1}c_1+c_2i = 
\sum_{i=1}^{n-1}c_1+\sum_{i=1}^{n-1}c_2i = 
c_1\cdot (n-1)+c_2\sum_{i=1}^{n-1}i = 
c_1\cdot (n-1)+c_2\frac{(n-1)n}{2} =\Theta(n^2).$$
The advantage of including the constants is that we can stop short of the final
step and get a better estimate of the actual number of operations used by the algorithm.
In other words, if we want an exact answer, we need to include the constants and lower
order terms.
If we just want a bound, the constants and lower order terms can often be ignored.
\end{exa}

\begin{rem}
There are rare cases when ignoring constants and lower order terms can cause 
trouble (meaning that it can lead to an incorrect answer) for subtle reasons
that are beyond the scope of this book. 
Unless you take more advanced courses dealing with these topics, 
you most likely won't run into those problems.
\end{rem}

Let's complicate things a bit by bringing in some basic data structures.

\begin{exa}\index{insertion sort}\index{sort!insertion}
Consider the following implementation of insertion sort that works on lists of integers
(written using Java syntax).
\begin{code}
void insertionSort(List<Integer> A, int n) {
    for (int i = 1; i < n; i++) {
        Integer T = A.get(i);
        int j = i-1;
        while (j >= 0 && A.get(j).compareTo(T) > 0) {
            A.set(j + 1, A.get(j));
            j--;
        }
        A.set(j+1, T);
    }
}
\end{code}
{\em Note: In some languages the size of the list, $n$, does not need to be passed in
since it can be obtained with a method call (e.g. \verb+A.size()+). 
We will pass it in just to be clear that the size of the list is $n$.}

In this implementation \verb+A.get(i)+ retrieves that $i$th element from the list
and \verb+A.set(i,x)+ sets the $i$th element of the list to $x$.
Also, 
 \verb+A.get(j).compareTo(T)+ returns a positive number if \verb+A.get(j)+
 is greater than \verb+T+, 0 of they are the same, and a negative number otherwise.
 It is essentially equivalent to \verb+A.get(j) > T+.

Give the complexity of this version of \verb+insertionSort+ assuming that the list is an array-based
implementation (e.g. an ArrayList in Java).  In case you are not aware from a previous course,
this means that both \verb+get+ and \verb+set+ take constant time.
In addition, we will assume that \verb+compareTo+ always takes constant time
(which should be true for any reasonable implementation).

\begin{sol}
Notice that this algorithm is almost identical to the earlier version that is implemented
on an array.  Array indexing and assignment are simply replaced with calls to \verb+get+
and \verb+set+.  Since the time of these calls remains constant, the earlier analysis still
holds.  Thus, the algorithm has a complexity of $\Theta(n^2)$.
\end{sol}
\end{exa}

The following should be review from a previous course. If you have not 
had a data structures course, read the solution and just take it as a given.

\begin{ques}
What are the complexities of the methods 
\verb+set(i,x)+ (set the $i$th element of the the list to $x$) and 
\verb+get(i)+ (return the $i$th element of the list) for a
{\em linked list}, assuming a reasonable implementation?

\noindent
\answerlinetwo
\begin{answer}
Since both of these methods require accessing the $i$th element of the list
for some integer $i$, and since we must traverse the list from the head, clearly
the complexity of both methods is $\Theta(i)$.
  
We could be less specific and say that the complexity is $\Theta(n)$ since $0\leq i <n$.  
However, when analyzing algorithms that make repeated calls to these methods,
using $\Theta(i)$ might give a more accurate answer overall.  It makes the analysis 
more difficult, but sometimes it is worth it.

{\em Note: For doubly-linked lists, 
some implementations traverse starting at the tail if the index is closer to the end
of the list.  However, that just means the complexity is no worse than $\Theta(n/2)=\Theta(n)$.
In other words, it only changes the complexity by a constant factor.}
\end{answer}
\end{ques}

\begin{exa}
Analyze the previous \verb+insertionSort+ algorithm assuming that the
list is now a {\em linked list}. 
\begin{sol}
The analysis here is a bit more complicated than we have previously seen, but we can still do it.
We start by analyzing the code inside the \verb+while+ loop.
In the worst case, each iteration of the loop makes two calls to \verb+A.get(j)+
and one call to \verb@A.set(j+1)@ and a constant amount of other work.  The total time for each iteration is therefore
about $2j+(j+1)+c=3j+c+1=3j+c^{\prime}$, where $c^{\prime}=c+1$ is still just some constant.  
The index of the while loop starts at $j=i$ and can go until
$j=1$ (with $j$ decrementing each iteration).  Thus, the complexity of the while loop
is about
$$\sum_{j=1}^{i}(3j+c^{\prime} )
= 3\sum_{j=1}^{i}j + \sum_{j=1}^{i} c^{\prime} 
=3\frac{(i+1)i}{2} +ic^{\prime} =\Theta(i^2)$$
The rest of the code inside the for loop takes constant time, except
the one call to \verb+get+ and one call to \verb+set+ which take time
$\Theta(i)$.\footnote{Note that there is a tricky part here.  It is subtle, but important.
The call to $set$ actually takes $j$ as a parameter.  However we cannot
use $\Theta(j)$ as the complexity of this because the $j$ has no meaning in this
context.  Therefore we use the fact that $1\leq j \leq i$ to instead call it
$\Theta(i)$.}
Thus, the code inside the for loop has complexity
$\Theta(i^2+i)=\Theta(i^2)$.
The outer for loop makes $i$  go from $1$ to $n-1$.
Thus, the overall complexity is 
$$\sum_{i=1}^{n-1}\Theta(i^2)
=\Theta\left(\sum_{i=1}^{n-1}i^2\right)
=\Theta\left(\frac{(n-1)n(2(n-1)+1)}{6}\right)
=\Theta(n^3).
$$
Clearly using a linked list in this implementation of insertion sort is a bad idea. 
\end{sol}
\end{exa}

You will get a chance to analyze the third popular basic sorting algorithm, {\em Selection Sort}, in the exercises.


%%%%%%%%%%%%%%%%%%%%%%%%%%%%%%%%%%%%%%%%%
\subsection{Basic Data Structures}\label{subsec:BDS}
If you have had a previous data structures course, this section should be review.
If you have not, you should probably skip this section since you will be unable to
complete the exercises without understanding the data structures involved.
%Now is a good time to review the complexities of operations on some
%common data structures (It is assumed that you have either learned this material before
%or that you can work them out based on your knowledge of how these data structures are implemented).

\begin{exer}\index{stack}
For each of the following implementations of a {\em stack}, 
give a tight bound (using $\Theta$-notation, of course) 
on the expected running time of the given operations, 
assuming that the data structure has $n$ items in it before the operation is performed.

\noindent
{\LARGE
\begin{tabular}{|l|l|l|}
\hline
{\bf Stack}&\, \, \,  array\, \, \, &\, \, linked list\, \,\\
\hline
\hline
push&&\\
\hline
pop&&\\
\hline
peek&&\\
\hline
size&&\\
%\hline
%isEmpty&&\\
\hline
\end{tabular}
}
\begin{answer}
All of them should be $\Theta(1)$, assuming we keep track of how many elements
are currently in the stack (which is a reasonable thing to do).
\end{answer}
\end{exer}

\begin{exer}\index{queue}
For each of the following implementations of a {\em queue}, 
give a tight bound 
on the expected running time of the given operations, 
assuming that the data structure has $n$ items in it before the operation is performed.

\noindent
{\LARGE
\begin{tabular}{|l|l|l|l|}
\hline
{\bf Queue}&\, \, \,  array\, \, \, &\, \, linked list\, \, &circular array\\
\hline
\hline
enqueue&&&\\
\hline
dequeue&&&\\
\hline
first&&&\\
\hline
size&&&\\
%\hline
%isEmpty&&&\\
\hline
\end{tabular}
}
\begin{answer}
For an array, either enqueue or dequeue (but not both) will be $\Theta(n)$.  
All of the others will be $\Theta(1)$, assuming we keep track of how many elements
are currently in the queue.  Note that the advantage of the circular array implementation
is that both enqueue and dequeue are $\Theta(1)$.
\end{answer}
\end{exer}


\begin{exer}
For each of the following implementations of a {\em list}, 
give a tight bound
on the expected running time of the given operations, 
assuming that the data structure has $n$ items in it before the operation is performed.

\noindent
{\LARGE
\begin{tabular}{|l|l|l|}
\hline
{\bf List}&\, \, \,  array\, \, \, &\, \, linked list\, \,\\
\hline
\hline
addToFront&&\\
\hline
addToEnd&&\\
\hline
removeFirst&&\\
\hline
contains&&\\
\hline
size&&\\
\hline
isEmpty&&\\
\hline
\end{tabular}
}
\begin{answer}
For the array implementation, addToFront, removeFirst and contains will all be $\Theta(n)$
and the rest will be $\Theta(1)$.
For the linked list implementation, contains will be $\Theta(n)$ and the rest will be 
constant if we assume there is both a head and tail pointer.  If there is no tail pointer,
addToEnd will be $\Theta(n)$.
\end{answer}
\end{exer}

\begin{exer}\index{binary search tree}\index{BST}
For each of the following implementations of a {\em binary search tree (BST)}, 
give a tight bound
on the expected running time of the given operations, 
assuming that the data structure has $n$ items in it before the operation is performed.
Assume a linked implementations (rather than arrays). 
For balanced, assume an implementation like red-black tree or AVL tree.

\noindent
{\LARGE
\begin{tabular}{|l|l|l|}
\hline
{\bf BST}&unbalanced&balanced\\
\hline
\hline
insert/add&&\\
\hline
delete/remove&&\\
\hline
search/contains&&\\
\hline
maximum&&\\
\hline
successor&&\\
\hline
\end{tabular}
}
\begin{answer}
For unbalanced, all operations can be done in $\Theta(h)$ time, where $h$ is the height 
of the tree.  This is the best answer you can give.  You could also say $O(n)$ time since
the height is no more then $n$, but this answer is not precise enough to be of much use.
You {\em cannot} say $\Theta(\log n)$ since 
this is not necessarily true for an unbalanced tree.

For balanced (red-black, AVL, etc.), all operations can be implemented with complexity $\Theta(\log n)$.
\end{answer}
\end{exer}


\begin{exer}\index{hash table}
Give the average- and worst-case complexity of the following operations
on a {\em hash table} (implemented with open-addressing or chaining--it doesn't matter),
assuming that the data structure has $n$ items in it before the operation is performed.

\noindent
{\LARGE
\begin{tabular}{|l|l|l|}
\hline
{\bf Hash Table}&average&\, \, worst\, \, \\
\hline
\hline
insert/add&&\\
\hline
delete/remove&&\\
\hline
search/contains&&\\
\hline
\end{tabular}
}
\begin{answer}
The average-case complexity for all of these operations is $\Theta(1)$, 
and the worst-case complexity is $\Theta(n)$.
\end{answer}
\end{exer}


\newpage
%%%%%%%%%%%%%%%%%%%%%%%%%%%%%%%%%%%%%%%%%
\subsection{More Examples}
Before presenting several more examples of algorithm analysis, let's
summarize a few principles from the examples we have seen so far.  
\begin{enumerate}
\item
We can usually replace constants with $1$.  For instance, if something performs 
$30$ operations, we can say it is constant and call it $1$.  This is only valid if 
it really is always $30$, of course.  
\item
We can usually ignore lower-order terms.
So if an algorithm takes $c_1n+c_2$ operations, we can usually say that it takes $n$.
\item
Nested loops must be treated with caution.  If the limits in an inner loop change based
on the outer loop, we generally need to write this as a summation.
\item
We should generally work from the inside-out.  Until you know how much time it takes
to execute the code inside a loop, you cannot determine how much time the loop takes.
\item
Function calls must be examined carefully.  Do not assume that a function
takes constant time unless you know that to be true. 
We already saw a few examples where function calls did {\em not}
take constant time, and the next example will demonstrate it again.
\item
Only the size of the input should appear as a variable in the complexity of an algorithm.
If you have variables like $i$, $j$, or $k$ in your complexity 
(because they were indexes of a loop, for instance), you should probably rethink
your analysis of the algorithm.  
Loop variables should {\em never} appear in the complexity of an algorithm.
\end{enumerate}

Now it's time to see if you can spot where someone didn't follow some of these principles.

\newpage

\begin{eval}\label{eval:addpow}
Consider the following code that computes $a^0+a^1+a^2+\cdots+a^{n-1}$.
\begin{code}
	  double addPowers(double a, int n) {
	      if(a==1) {
	          return n;
	      } else {
	          double sum = 0;
	          for(int i=0;i<n;i++) {
	              sum += power(a,i); 
	          }
	          return sum;
	      }
	  }
\end{code}  
The function \verb+power(a,i)+ computes $a^i$, and takes $i$ operations.
Regard the input size as $n$.
What is the worst-case complexity of \verb+addPowers(a,n)+?
\begin{ssolenum}
\item
Since $a^n$ is an exponential function, the complexity is $O(a^n)$.

\noindent
\evallinethree
\item
The worst-case is $n i$ since \verb+power(a,i)+ takes $i$ time and the for loop
executes $n$ times.

\noindent
\evallinethree
\item
The for loop executes $n$ times.  Each time it executes, it calls \verb+power(a,i)+,
which takes $i$ time.  In the worst case, $i=n-1$, so the complexity is
$(n-1)n=O(n^2)$.

\noindent
\evallinethree
\end{ssolenum}
\begin{answer}
\begin{solevalenum}
\item
I have no idea what logic they are trying to use here.  Sure, $a^n$ is an exponential
function, but what does that have to do with how long this algorithm takes?
This solution is way off.
\item
Having the $i$ in the answer is nonsense since it doesn't mean
anything in the context of a complexity---it is just a variable that happens to be
used to index a loop.  Further, the answer should be given using
$\Theta$-notation.  So this solution is just plain wrong.
Since having an $i$ in the complexity does not make any sense, this person either
has a fundamental misunderstanding of how to analyze algorithms or they didn't 
think about their final answer.  Don't be like this person!
\item
This solution is O.K., but it has a slight problem.  Although the analysis given estimates
the worst-case behavior, it over-estimates it.  By replacing $i$ with $n-1$, they are
over-estimating how long the algorithm takes. 
The call to \verb+pow+ only takes $n-1$ time once.
This solution can really only tell us
that the complexity is $O(n^2)$.  Is it possible the over-estimation of time resulted
in a bound that isn't tight?  Even if it turns out that $\Theta(n^2)$ is the correct bound,
this solution does not prove it.
Although they are on the right track, 
this person needed to be a little more careful in their analysis.
\end{solevalenum}
\end{answer}
\end{eval}

\newpage

\begin{exer}
What is the worst-case complexity of \verb+addPowers+ from Evaluate~\ref{eval:addpow}?
Justify your answer.

\vspace*{2.25in}
\begin{answer}
Before you read too far: 
if you did not use a summation in your solution, go back and try again!
This is very similar to the analysis of \verb+bubblesort+.  The for loop takes $i$ from
$0$ to $n-1$, and each time the code in the loop takes $i$ time (since that is how long
\verb+power(a,i)+ takes).  Thus, the complexity is 
$$\sum_{i=0}^{n-1} i = \frac{(n-1)n}{2}=\Theta(n^2).$$
Notice that just because the answer is $\Theta(n^2)$, that does not mean that the
third solution to Evaluate~\ref{eval:addpow} was correct.  As we stated in the solution
to that problem, because they overestimated the number of operations,
they only proved that the algorithm has complexity $O(n^2)$.
\end{answer}
\end{exer}

\begin{exer}
Give an implementation of the \verb+addPowers+ algorithm that takes $\Theta(n)$ time.
Justify the fact that it takes $\Theta(n)$ time. 
(Hint: Why compute $a^5$ (for instance) from scratch if you have already computed $a^4$?)
\begin{code}
	  double addPowers(double a, int n) {
	     
	     
	     
	     
	     
	     
	     
	     
	     
	     
	     
	     
	     
	     
	     
	     
	  }
\end{code}
Justification of complexity:

\vspace*{1.25in}
\begin{answer}
Here is one solution.
\begin{code}
	  double addPowers(double a, int n) {
	      if(a==1) {
	          return n;
	      } else {
	          double sum = 1; // for the $a^0$ term.
	          double pow = 1;
	          for(int i=1;i<n;i++) {
	              pow = pow*a;
	              sum += pow; 
	          }
	          return sum;
	      }
	  }
\end{code} 
If $a=1$, the algorithm takes constant time.
Otherwise, it executes a constant number of operations and a single for loop $n$ times.
The code in the loop takes constant time.  
Thus the algorithm takes $\Theta(n)$ time.
\end{answer} 
\end{exer}

\newpage % another awkward break

\begin{exer}
Give an implementation of the \verb+addPowers+ algorithm that takes $\Theta(n)$ time
{\em but does not use a loop}.  Justify the fact that it takes $\Theta(n)$ time.
(Hint:  This solution should be much shorter than your previous one.)
\begin{code}
	  double addPowers(double a, int n) {
	     
	     
	     
	     
	     
	     
	     
	     
	     
	     
	     
	     
	     
	     
	  }
\end{code}
Justification of complexity:

\vspace*{1.25in}
\begin{answer}
If you used recursion instead of a loop, cool idea.  However, go back and do it again.
There is an even {\em simpler} way to do it.  Need a hint?  Apply some of that
discrete mathematics material you have been learning!  When you have a solution that
does not use a loop or recursion (or you get stuck), keep reading.

The trick is to use the formula for a
geometric series (did you recognize that this is what \verb+addPowers+ is really computing?).  
We need a special case for $a=1$ because the 
formula requires that $a\neq 1$. 

\begin{code}
	  double addPowers(double a, int n) {
	      if(a==1) {
	          return n;
	      } else {
	          return (1-power(a,n))/(1-a);
	      }
	  }
\end{code} 

If $a=1$, the algorithm takes constant time.
Otherwise, it executes a constant number of operations and a single call
to \verb@power(a,n)@ which takes $n$ time.
Thus the algorithm takes $\Theta(n+1)=\Theta(n)$ time.

It is worth noting that $a=0$ is a tricky case.
\verb+addPowers+ can't really be computed for $a=0$ since $0^0$ is undefined.
It is for this reason that the first term of a geometric sequence is technically $1$, not $a^0$.
Since $a^0=1$ for all other values of $a$, the case of $a=0$ is usually glossed over.
If you don't understand what the fuss is about, don't worry too much about it.
\end{answer} 
\end{exer}

\begin{exa}\label{exa:summn}
A student turned in the code below (which does as its name suggests).  
I gave them a `C' on the assignment because although it works,
it is very inefficient.  About how many operations does their
implementation require?
\begin{code}
	int sumFromMToN(int m, int n) {
	    int sum = 0;
	    for(int i=1;i<=n;i++) {
	        sum = sum + i;
	    }
	    for(int i=1;i<m;i++) {
	        sum = sum - i;
	    }
	    return sum;
	}
\end{code}
\begin{sol}
The first loop takes about $1+4n$ operations, and the second loop takes about
$1+4(m-1)$ operations.  The first statement and return statement add 2 operations.
So the total number of operations is about $4+4n+4(m-1)=4(n+m)=\Theta(n+m)$.
\end{sol}
\end{exa}


\begin{eval}
Write an `A' version of the method from Example~\ref{exa:summn}.
You can assume that $1\leq m\leq n$. 
{\em For each solution, determine how many operations are required and
evaluate it based on that as well as whether
or not it is correct.}
\begin{ssolenum}
\item 
\begin{code}

	int sumFromMToN(int m,int n) {
	    int sum = 0;
	    for(int i=0;i<n;i++) {
	        sum = sum + i;
	    }
	    for(int i=0;i<m;i++) {
	        sum = sum - i;
	    }
	    return sum;
	}
\end{code}
\evallinethree

\item
\begin{code}

	int sumFromMToN(int m,int n) {
	    int sum = 0;
	    for(int i=m;i<n;i++) {
	        sum = sum + i;
	    }
	    return sum;
	}
\end{code}
\evallinethree

\item
\begin{code}

	int sumFromMToN(int m,int n) {
	    return (n*(n-1)/2 - m*(m-1)/2);
	}
\end{code}
\evallinethree

\end{ssolenum}
\begin{answer}
\begin{solevalenum}
\item
This takes about $2+4n+4m$ operations, which is essentially the
same as the `C' version.  Unfortunately, it is slightly worse than the original
solution since it is now incorrect.  All they did is omit adding the final $n$ in the first
sum.  This went from a `C' to a `D' (at best).
\item
This one takes about $2+4(n-1-m+1)=2+4(n-m)$ operations.  This is a lot better than
the previous solutions.  Unfortunately, it misses adding the final $n$, so it is incorrect.
It also is not as efficient as possible.  I'd say this is still a `C'.
\item
This student figured out the trick---they know a formula to compute the sum, 
so they tried to use it.  
Unfortunately, they used the formula incorrectly and/or they made a mistake when
manipulating the sum (it is impossible to tell exactly what they did wrong---they 
made either one or two errors), so the algorithm is not correct.
In terms of efficiency, their solution is great because it takes a constant
number of operations no matter what $n$ and $m$ are.    
Because their answer is efficient and {\em very close} to being correct, 
I'd probably give them a `B'.
\end{solevalenum}
\end{answer}
\end{eval}

\begin{exer}
Write an `A' version of the method from Example~\ref{exa:summn}.
You can assume that $1\leq m\leq n$. 
Explain why your solution is correct and give its efficiency.
\begin{code}
	int sumFromMToN(int m,int n) {
	    
	    
	    
	    
	    
	    
	    
	    
	}
\end{code}

\noindent
Justification

\vspace*{1in}
\noindent
Efficiency with justification

\vspace*{1in}
\begin{answer}
We use the fact that 
${\displaystyle \sum_{k=m}^n k = \sum_{k=1}^n - \sum_{k=1}^{m-1}
=\frac{n(n+1)}{2} - \frac{(m-1)m}{2} }$ to give us the following solution:
\begin{code}
    int sumFromMToN(int m,int n) {
	    return n*(n+1)/2 - (m-1)*m/2;
	}
\end{code}
Since this is just doing a fixed number of arithmetic operations no matter
what the values of $m$ and $n$ are, it takes constant time.
\end{answer}
\end{exer}

\begin{exa}\label{exa:matmult}\index{matrix multiplication}
\index{matrix!multiplication}
The \verb+MatrixMultiply+ algorithm given below is the standard algorithm used to 
compute the product of two matrices.
Find the worst-case complexity of \verb+MatrixMultiply+.
Assume that $A$ and $B$ are $n\times n$ matrices.
\begin{code}
  Matrix MatrixMultiply(Matrix A, Matrix B) {
      Matrix C;
      for(int i=0 ; i < n; i++) {
          for(int j=0 ; j < n ; j++) {
              C[i][j]=0;
              for(int k=0 ; k < n ; k++) {
                  C[i][j] += A[i][k]*B[k][j];
              }
          }
      }
      return C;
  }
\end{code}
\begin{sol}
The code inside the inner loop does array indexing, multiplication, addition,
and assignment.  All of these together take just constant time.  
Therefore, let's count the number of times the statement \verb#C[i][j]+=A[i][k]*B[k][j]#
executes.  We will ignore the calls to \verb+C[i][j]=0+ since it executes just once
every time the entire middle loop executes, so it has a negligible contribution.
Similarly, the statement \verb#C[i][j]+=A[i][k]*B[k][j]# is called at least as often
as any of the code in the for loops (i.e. the comparisons and increments) so we will
ignore that code as well.  The bottom line is that if we count the number of times 
\verb#C[i][j]+=A[i][k]*B[k][j]# executes, it will give us a tight bound on the complexity
of \verb+MatrixMultiply+.

The inner loop executes the statement $n$ times.  
The middle loop executes $n$ times, each time executing the inner loop (which executes
the statement $n$ times).
Thus, the middle loop executes the statement $n\times n=n^2$ times.
The outer loop simply executes the middle loop $n$ times.  Therefore the outer loop
(and thus the whole algorithm) executes the statement $n\times n^2=n^3$ times.
Thus, the worst-case complexity of \verb+MatrixMultiply+ is $\Theta(n^3)$.
Notice that this is also the best and average-case complexity since there are no
conditional statements in this code.
\end{sol}
\end{exa}

\begin{exa}\label{exa:retainAllAL}
In Java, the ArrayList \verb+retainAll+ method is implemented as follows 
(this code is simplified a little from the actual implementation, but the changes do not affect the complexity of the code).  
Note that \verb+Object[] elementData+ and \verb+int size+ are fields of ArrayList whose meaning should be obvious.
\begin{code}
  public boolean retainAll(Collection<?> c) {
      boolean modified = false;
      int w = 0;
      for (int r = 0; r < size; r++) {
          if (c.contains(elementData[r])) {
              elementData[w++] = elementData[r];
          }
      }
      if (w != size) {
          for (int i = w; i < size; i++)
              elementData[i] = null;
          size = w;
          modified = true;
      }
      return modified;
  }
\end{code}
Let \verb+al1+ be an ArrayList of size $n$.
\begin{lenum}
\item
What is the complexity of \verb+al1.retainAll(al2)+, where \verb+al2+ is an ArrayList with $m$ elements?
\begin{sol}
The method call \verb+c.contains(elementData[r])+ takes $\Theta(m)$ time since $c$
is an ArrayList.  The rest of the code in that for loop takes constant time.
Since this is done inside a for loop that executes $n$ times, the first 
half of the code takes $\Theta(n m)$ time.
In the worst case ($w=0$), the second half of the code takes $\Theta(n)$ time.
Thus, the worst-case complexity of the method is $\Theta(n m+n)=\Theta(n m)$.
\end{sol}
\item
What is the complexity of \verb+al1.retainAll(ts2)+, where \verb+ts2+ is a TreeSet with $m$ elements?
\begin{sol}
The method call \verb+c.contains(elementData[r])+ takes $\Theta(\log m)$ time since $c$
is a TreeSet.  The rest of the code in that for loop takes constant time.
Since this is done inside a for loop that executes $n$ times, the first 
half of the code takes $\Theta(n \log m)$ time.
In the worst case ($w=0$), the second half of the code takes $\Theta(n)$ time.
Thus, the worst-case complexity of the method is $\Theta(n \log m+n)=\Theta(n \log m)$.
\end{sol}
\end{lenum}
\end{exa}

\begin{exer}
Answer the following two questions based on the code from Example~\ref{exa:retainAllAL}.
\begin{lenum}
\item
What is the complexity of \verb+al1.retainAll(ll2)+, 
where \verb+ll2+ is a LinkedList with $m$ elements?
\answerlinetwo
\item
What is the complexity of \verb+al1.retainAll(hs2)+, 
where \verb+hs2+ is a HashSet with $m$ elements?
\answerlinetwo
\end{lenum}
\begin{answer}
The analysis of these is very similar to the analysis of Examples~\ref{exa:retainAllAL}, so the details are omitted.
(a) For a LinkedList the \verb+contains+ method takes $\Theta(m)$ time, so the overall
complexity is $\Theta(nm)$.
(b)For a HashSet it takes $\Theta(1)$ time to call \verb+contains+ (on average), so
the overall complexity is $\Theta(n+n)=\Theta(n)$.
\end{answer}
\end{exer}








\begin{exa}\label{exa:retainAllOthers}
In Java, the \verb+retainAll+ method is implemented as follows for
\verb+LinkedList+, \verb+TreeSet+, and \verb+HashSet+.
\begin{code}
    public boolean retainAll(Collection<?> c) {
        boolean modified = false;
        Iterator<E> iter = iterator();
        while (iter.hasNext()) {
            if (!c.contains(iter.next())) {
                iter.remove();
                modified = true;
            }
        }
        return modified;
    }
\end{code}
Assume that the calls \verb+iter.hasNext()+ and \verb+iter.next()+
take constant time.  
Let \verb+ts1+ be a TreeSet of size $n$. 
Find the worst-case complexity of each of the following method calls.

\begin{lenum}
\item
\verb+ts1.retainAll(al2)+, where \verb+al2+ is an ArrayList of size $m$.
\begin{sol}
The call to \verb+iter.remove()+ takes $\Theta(\log n)$ time since the iterator
is over a TreeSet.  The call to \verb+contains+ takes $\Theta(m)$ time since in this
case $c$ is the ArrayList \verb+al2+.   The other operations in the loop take constant time.
Thus, each iteration of the while loop takes $\Theta(\log n+m)$ time in the worst case
(which occurs if the conditional statement is always true and \verb+remove+ is called
every time). 
Since the
loop executes $n$ times, and the rest of the code takes constant time, 
the overall complexity is $\Theta(n(\log n + m))$.
\end{sol}
\item
\verb+ts1.retainAll(hs2)+, where \verb+hs2+ is an HashSet of size $m$.
\begin{sol}
The call to \verb+iter.remove()+ takes $\Theta(\log n)$ time.  
The call to \verb+contains+ takes $\Theta(1)$ time since $c$ is the HashSet \verb+hs2+.  
Thus, each iteration of the while loop takes $\Theta(\log n+1)$ time and the
overall complexity is therefore $\Theta(n(\log n + 1))=\Theta(n\log n)$.
\end{sol}
\end{lenum}
\end{exa}

\begin{exer}
Using the setup and code from Example~\ref{exa:retainAllOthers}, 
determine the complexity of the following method calls.
\begin{lenum}
\item
\verb+ts1.retainAll(ll2)+, where \verb+ll2+ is a LinkedList of size $m$.
\answerlinethree
\item
\verb+ts1.retainAll(ts2)+, where \verb+ts2+ is an TreeSet of size $m$.
\answerlinethree
\end{lenum}
\begin{answer}
(a) contains takes $\Theta(m)$ time so the complexity is
$\Theta(n(\log n+m))$.

(b) Here the contains method takes $\Theta(\log m)$ time, so the overall
complexity is $\Theta(n(\log n+\log m))$ 
(or $\Theta(n\log n + n\log m)$
if you prefer to write it that way).

Note that we don't know the relationship between
$n$ and $m$ so we can't simplify either answer.
\end{answer}
\end{exer}

It is important to note that the number of examples related to the \verb+retainAll+ method
is not reflective of the importance of this method.  It just turns out to be an 
interesting method to analyze the complexity of given different data structures.


We end this section with a comment that perhaps too few people think about.
{\em Theory} and {\em practice} don't always agree. 
Since asymptotic notation ignores the {\em constants}, two algorithms that have the same
complexity are not always equally good in practice.  
For instance, if one takes $4\cdot n^2$ operations
and the other $10,000\cdot n^2$ operations, clearly the first will be preferred even though
they are both $\Theta(n^2)$ algorithms.

As another example, consider matrix multiplication, which is used extensively in
many scientific applications.
As we saw, the standard algorithm has complexity $\Theta(n^3)$.  Strassen's 
algorithm for matrix multiplication (the details of which are beyond the scope of this book)
has complexity of about $\Theta(n^{2.8})$.
Clearly, Strassen's algorithm is better asymptotically.  In other words, if your matrices
are large enough, Strassen's algorithm is certainly the better choice.  However, it turns
out that if $n=50$, the standard algorithm performs better.
There is debate about the ``crossover point.''  This is the point at which the more efficient 
algorithm is worth using.  For smaller inputs, the overhead associated with the cleverness
of the algorithm isn't worth the extra time it takes.  For larger inputs, the extra overhead
is far outweighed by the benefits of the algorithm.  For Strassen's algorithm, this point
may be somewhere between $75$ and $100$, but don't quote me on that.
The point is that for small enough matrices, the standard algorithm should be used.
For matrices that are large enough, Strassen's algorithm should be used.
Neither one is {\em always} better to use.

Analyzing recursive algorithms can be a little more complex.  
We will consider such algorithms in Chapter \ref{ch:recurrences}, 
where we develop the necessary tools.

\subsection{Binary Search}\label{section:binarySearchAnalysis}
\index{binary search}
\index{search!binary}
Next we analyze the {\em binary search} algorithm.  
If you haven't seen it before, the concept is simple: 
If you want to find something in a sorted list, start by comparing with the middle element.
If the array is empty, the element is not in it.
If they are the middle and the element you are searching for are the same, you have found what you are looking for.
If what you are looking for is smaller than the middle element, then it is in the first half of the array.
If what you are looking for is bigger than the middle element, then it is in the second half of the array.
Now we simply repeat the process on the appropriate half of the array until we have our answer.

\begin{exa}
Here is an implementation of {\em binary search}.
\begin{code}
    int binarySearch(int a[], int n, int val) {
       int left=0, right=n-1;
       while (right-left>=0) {
           int middle = (left+right)/2;
           if(val==a[middle]) 
               return middle;
           else if(val<a[middle]) 
               right=middle-1;
           else 
               left=middle+1; 
        }
       return -1;
    }
\end{code}
\end{exa}

You may be familiar with the
recursive version of this algorithm instead of this iterative implementation.
In some ways it is more intuitive, but we have not yet covered recursion
or the analysis of recursive algorithms so instead we will analyze the iterative version.

Before we can analyze the algorithm, we need to develop a few useful results that will make the
proof much easier to understand.
We start by trying to get you to understand how the binary representation of $n$ and
$\lfloor n/2\rfloor$ are related to each other.

\begin{exa}
How is the binary representation of a number $n$ related to the
binary representation of $\lfloor n/2\rfloor$?  Let's try some examples.
If $n=9$, $\lfloor n/2\rfloor=4$.  Notice that the binary representation of $9$ is $1001$
and the binary representation of $4$ is $100$.  
If $n=22$, $\lfloor n/2\rfloor=11$.  The binary representation of $22$ is $10110$
and the binary representation of $11$ is $1011$.  
Is there a pattern here?  This probably isn't enough data to be certain yet.  
\end{exa}

Let's see if you can find the pattern with just a few more data points.

\begin{exer}
Fill in the following table with the binary representations.
\begin{center}
{\Large
\begin{tabular}{|c|c|c|c|}
\hline
\multicolumn{2}{|c|}{$n$}&
\multicolumn{2}{c|}{$\lfloor n/2\rfloor$}\\
decimal&\, \, \, binary\, \, \, &decimal&\, \, \, binary\, \, \, \\
\hline
\hline
12&&6&\\
\hline
13&&6&\\
\hline
32&&16&\\
\hline
33&&16&\\
\hline
118&&59&\\
\hline
119&&59&\\
\hline
\end{tabular}
}
\end{center}

\begin{answer}
\begin{tabular}{|l|l|l|l|l|}
\hline
\multicolumn{2}{|c|}{$n$}&
\multicolumn{2}{c|}{$\lfloor n/2\rfloor$}\\
decimal&binary\, \, \, \, \, &decimal&binary\, \, \, \, \,\\
\hline
\hline
12&1100&6&110\\
\hline
13&1101&6&110\\
\hline
32&100000&16&10000\\
\hline
33&100001&16&10000\\
\hline
118&1110110&59&111011\\
\hline
119&1110111&59&111011\\
\hline
\end{tabular}
\end{answer}
\end{exer}

\begin{ques}
How are the binary representations of 
$n$ and $\lfloor n/2\rfloor$ related?

\noindent
\answerlinetwo
\begin{answer}
The next theorem answers the question about the pattern.
\end{answer}
\end{ques}

Hopefully you observed a clear pattern in the previous exercise.  The next theorem formalizes 
this idea.  We provide a proof of the theorem to make it clear what is going on.

\begin{thm}\label{thm:binnovertwo}
The binary representation of $\lfloor n/2\rfloor$ 
is the binary representation of n shifted to the right one bit.  
That is, the binary representation of $\lfloor n/2\rfloor$  is the same as that of 
$n$ with the last bit (the lowest order bit) chopped off.
\begin{pf}
Let the binary representation of
$n$ be $a_m a_{m-1}a_{m-2}\ldots a_2a_1a_0$, where $a_m=1$
(so the highest order bit is a $1$).  
Then 
$$n=a_m\, 2^m + a_{m-1}\, 2^{m-1}+ \ldots+ a_2\, 2^2 + a_1\, 2^1 +a_0\, 2^0.$$
From this we can see that
\begin{eqnarray*}
\lfloor n/2\rfloor
&=&\lfloor (a_m\, 2^m + a_{m-1}\, 2^{m-1} + 
    \ldots+ a_2\, 2^2 + a_1\, 2^1 + a_0\, 2^0)/2\rfloor \\
&=&  \lfloor a_m\, 2^m/2 + a_{m-1}\, 2^{m-1}/2 + 
    \ldots+ a_2\, 2^2/2 + a_1\, 2^1/2 + a_0\, 2^0/2\rfloor \\
&=&  \lfloor a_m\, 2^{m-1} + a_{m-1}\, 2^{m-2} + 
    \ldots+ a_2\, 2^1 + a_1\, 2^0 + a_0/2\rfloor\\
&=&  a_m\, 2^{m-1} + a_{m-1}\, 2^{m-2} + 
    \ldots+ a_2\, 2^1 + a_1\, 2^0.
\end{eqnarray*}
Notice that in the last step, $a_0/2$ is chopped off by the floor since it is
either $0/2$ or $1/2$ and the other numbers are integers.
From this we can see that the binary representation of 
$\lfloor n/2\rfloor$  is $a_ma_{m-1}a_{m-2}\ldots a_2a_1$, 
which is the binary representation of $n$ shifted to the right one bit.
\end{pf}
\end{thm}
\begin{cor}\label{cor:kbits}
If the number $n$ requires exactly $k$ bits to represent in binary, 
then $\lfloor n/2\rfloor$ requires exactly $k-1$ bits to represent in binary.
\begin{pf}
According to Theorem~\ref{thm:binnovertwo}, 
the binary representation of $\lfloor n/2\rfloor$ is the binary representation of
$n$ shifted to the right one bit.  
Thus it is clear that $\lfloor n/2\rfloor$ requires one less bit to represent.
\end{pf}
\end{cor}

We need just one more result.

\begin{thm}\label{thm:numbinarybits}
It takes $\lfloor \log_2 n\rfloor +1$ bits to represent $n$ in binary.
\begin{pf}
Recall that $\log_c b$ is defined as 
``the number that $c$ must be raised to in order to get $b$.'' 
That is, if $k = \log_c b$, then $c^k=b$.  
Also, it should be clear that $2^k$ is the smallest number that requires 
$k+1$ bits to represent in binary. (If you are not convinced of this, write out
some binary numbers near powers of two until you see it.)
Let $k$ be the number such that 
\begin{eqnarray}
\label{eqn:binequ}&& 2^{k-1}\leq n < 2^k.
\end{eqnarray}
Since writing $2^{k-1}$ takes $k$ bits and $2^k$ is the smallest number
that requires $k+1$ bits, it should be clear that 
$n$ requires exactly $k$ bits to represent in binary.  
Taking the logarithm of equation~\ref{eqn:binequ}, we get 
$$\log_2 2^{k-1}\leq \log_2 n < \log_2 2^k,$$ 
which leads to  
$$k-1\leq \log_2 n < k.$$
Clearly $\lfloor \log_2 n\rfloor =k-1$ since it is an integer.
Thus, 
$k=\lfloor \log_2 n\rfloor+1$, so it
takes $\lfloor \log_2 n\rfloor +1$ bits to represent $n$ in binary.
\end{pf}
\end{thm}

Now we are ready to analyze {\em binary search}.

\begin{exa}\label{exa:binSearchIter}\index{binary search}
We will show that {\em binary search} has worst-case complexity $\Theta(\log n)$.
More precisely, 
we will prove that the while loop executes no more than $\lfloor \log_2 n\rfloor +1$ times.
Here is the algorithm as a reminder:
\begin{code}
    int binarySearch(int a[], int n, int val) {
       int left=0, right=n-1;
       while (right-left>=0) {
           int middle = (left+right)/2;
           if(val==a[middle]) 
               return middle;
           else if(val<a[middle]) 
               right=middle-1;
           else 
               left=middle+1; 
        }
       return -1;
    }
\end{code}

Recall that binary search finds the index of a value in a sorted array
by comparing the value being searched for with the middle element of the array.
If they are the same, it returns the index of the element.  Otherwise it
continues the search in only half of the array.  In other words, it removes
from consideration half of the array at each iteration.  Which half depends on whether
the search value was greater than or less than the middle value.

Since the code inside the while loop takes a constant amount of time, the
complexity of {\em binary search} depends only on the number of iterations
of the loop.  Clearly the worst case is when a value is not in the array since
otherwise the loop ends early with the \verb+return+ statement.  Thus we will
assume the value is not in the array.

Notice that the value \verb+right-left+ is the number of entries of the array
that are still under consideration by the algorithm.
The loop executes until \verb+right-left+ $<0$.
Before the first iteration, \verb+right-left+ $=n$.
During each iteration,  
either \verb+right+ or \verb+left+ is set to the
middle value between \verb+right+ and \verb+left+ (plus or minus 1).
So after the first iteration, \verb+right-left+ $\leq \lfloor n/2\rfloor$.
In other words, the algorithm has discarded at least half of the entries of
the array.
During each subsequent iteration, \verb+right-left+ continues to be no more than the floor
of half of its previous value, so the algorithm continues to discard half of the entries
of the array each time through the loop.

According to Corollary~\ref{cor:kbits}, each iteration of the loop reduces 
the number of bits used to represent \verb+right-left+ by one.
According to Theorem~\ref{thm:numbinarybits}, 
it takes $\lfloor \log_2 n\rfloor +1$ bits to represent $n$ in binary, and
$right-left$ started out as $n$.
Therefore, after $\lfloor \log_2 n\rfloor$ iterations through the loop, \verb+right-left+
becomes $1$, and the next iterations ensures that \verb+right-left+ becomes
negative and the loop terminates (check this!).  
Since the loop executes at most $\lfloor \log_2 n\rfloor +1$  times,
the worst-case complexity of {\em binary search} is $\Theta(\log n)$.
\end{exa}


%%%%%%%%%%%%%%%%%%%%%%%%%%%%%%%%%%%%%%%%%%%%%%%%%%%%%%%%


